\documentclass[11pt,a4paper]{report}
\usepackage{../common/preamble}

% ---------- Pacchetti usati dal codice generato da pandoc ----------
\usepackage{longtable,array,calc}
\providecommand{\tightlist}{\setlength{\itemsep}{0pt}\setlength{\parskip}{0pt}}
\makeatletter\renewcommand{\verbatim@font}{\ttfamily\footnotesize}\makeatother

% ---------- Testi in italiano (babel italiano non disponibile ovunque) ----------
\addto\captionsenglish{%
  \renewcommand{\contentsname}{Indice}%
  \renewcommand{\chaptername}{Capitolo}%
  \renewcommand{\appendixname}{Appendice}%
  \renewcommand{\proofname}{Dimostrazione}%
}

% ---------- Caratteri Unicode presenti negli appunti ----------
\usepackage{newunicodechar}
\newunicodechar{ε}{\ensuremath{\varepsilon}}
\newunicodechar{λ}{\ensuremath{\lambda}}
\newunicodechar{τ}{\ensuremath{\tau}}
\newunicodechar{Π}{\ensuremath{\Pi}}
\newunicodechar{ℓ}{\ensuremath{\ell}}
\newunicodechar{←}{\ifdim\fontdimen3\font=0pt\makebox[0.525em]{\ensuremath{\leftarrow}}\else\ensuremath{\leftarrow}\fi}
\newunicodechar{→}{\ifdim\fontdimen3\font=0pt\makebox[0.525em]{\ensuremath{\rightarrow}}\else\ensuremath{\rightarrow}\fi}
\newunicodechar{↓}{\ifdim\fontdimen3\font=0pt\makebox[0.525em]{\ensuremath{\downarrow}}\else\ensuremath{\downarrow}\fi}
\newunicodechar{⇒}{\ensuremath{\Rightarrow}}
\newunicodechar{∈}{\ensuremath{\in}}
\newunicodechar{∎}{\ensuremath{\blacksquare}}
\newunicodechar{−}{\ensuremath{-}}
\newunicodechar{≈}{\ensuremath{\approx}}
\newunicodechar{≠}{\ensuremath{\neq}}
\newunicodechar{≥}{\ensuremath{\geq}}
\newunicodechar{⊕}{\ensuremath{\oplus}}
\newunicodechar{✓}{\ensuremath{\checkmark}}
\newunicodechar{€}{EUR}
\newunicodechar{─}{\makebox[0.525em]{\rule[0.55ex]{0.525em}{0.4pt}}}
\newunicodechar{│}{\makebox[0.525em]{\rule[-0.35ex]{0.4pt}{2.4ex}}}
\newunicodechar{┌}{\makebox[0.525em][r]{\rule[-0.35ex]{0.4pt}{1.3ex}\rule[0.55ex]{0.26em}{0.4pt}}}
\newunicodechar{┐}{\makebox[0.525em][l]{\rule[0.55ex]{0.26em}{0.4pt}\rule[-0.35ex]{0.4pt}{1.3ex}}}
\newunicodechar{└}{\makebox[0.525em][r]{\rule[0.55ex]{0.4pt}{1.5ex}\rule[0.55ex]{0.26em}{0.4pt}}}
\newunicodechar{┘}{\makebox[0.525em][l]{\rule[0.55ex]{0.26em}{0.4pt}\rule[0.55ex]{0.4pt}{1.5ex}}}

% ---------- Riquadri ----------
% inbreve: la definizione corta prima di ogni concetto importante
\newtcolorbox{inbreve}{enhanced, breakable,
  colback=RoyalBlue!6, colframe=RoyalBlue!60, boxrule=0.6pt, arc=2pt,
  left=6pt, right=6pt, top=8pt, bottom=6pt, before skip=12pt, after skip=10pt,
  fonttitle=\bfseries\color{black},
  attach boxed title to top left={xshift=8mm, yshift*=-\tcboxedtitleheight/2},
  boxed title style={colback=RoyalBlue!30, colframe=RoyalBlue!30, sharp corners, boxrule=0pt},
  title={In breve}}
\newtcolorbox{esempio}{enhanced, breakable,
  colback=TealBlue!5, colframe=TealBlue!55, boxrule=0.6pt, arc=2pt,
  left=6pt, right=6pt, top=6pt, bottom=6pt, before skip=10pt, after skip=10pt}
\newtcolorbox{nota}{enhanced, breakable,
  colback=Gray!6, colframe=Gray!45, boxrule=0.6pt, arc=2pt,
  left=6pt, right=6pt, top=6pt, bottom=6pt, before skip=10pt, after skip=10pt}
\newtcolorbox{approfondimento}[1]{enhanced, breakable,
  colback=white, colframe=Goldenrod!70, boxrule=0.6pt, arc=2pt,
  left=6pt, right=6pt, top=8pt, bottom=6pt, before skip=12pt, after skip=10pt,
  fonttitle=\bfseries\color{black},
  attach boxed title to top left={xshift=8mm, yshift*=-\tcboxedtitleheight/2},
  boxed title style={colback=Goldenrod!35, colframe=Goldenrod!35, sharp corners, boxrule=0pt},
  title={#1}}

\title{\textbf{Crittografia}\\[0.5em]\large Appunti di Luiggi\\[0.3em]\normalsize Prof.\ Daniele Venturi}
\author{Luiggi}
\date{Anno accademico 2026/2027}

\begin{document}
\maketitle

\noindent Gli appunti seguono quelli scritti a mano dal prof. Prima di ogni concetto importante c'è un riquadro \textbf{In breve} con la definizione corta. Le parti aggiunte per capire meglio sono segnate \textit{(non negli appunti del prof)}.

\tableofcontents

% Generato da tools/notion2latex.py a partire dalla pagina Notion "lezione-1".

\chapter{Lezione 1 -- Introduzione e segretezza perfetta}

\section{Info sul corso}

\begin{itemize}
\tightlist
\item
  \textbf{Docente}: Prof.~Daniele Venturi
\item
  \textbf{Esame}: scritto, 3 parti. Ogni parte ha \textbf{1 esercizio + 1 domanda di teoria}
\item
  \textbf{Libro di testo}: \emph{Introduction to Modern Cryptography}, J. Katz e Y. Lindell
\end{itemize}

\section{Di cosa parla il corso}

Il corso copre le basi della crittografia con un focus sulla \textbf{sicurezza dimostrabile} (\emph{provable security}), che è l'approccio moderno.

\begin{itemize}
\tightlist
\item
  \textbf{Prima del 1950 circa}: la crittografia era \emph{trial and error}. Si inventava un cifrario, qualcuno lo rompeva, se ne inventava un altro.
\item
  \textbf{Dal 1950 in poi}: diventa una \textbf{scienza}.
\end{itemize}

\textbf{Sicurezza dimostrabile} significa due cose:

\begin{enumerate}
\def\labelenumi{\arabic{enumi}.}
\tightlist
\item
  dare una \textbf{definizione precisa} di cosa vuol dire ``sicuro''
\item
  fornire una \textbf{dimostrazione} che lo schema soddisfa quella definizione (eventualmente sotto alcune ipotesi)
\end{enumerate}

\subsection{Obiettivo: comunicazione sicura}

Scenario classico: \textbf{Alice} manda un messaggio \(m\) a \textbf{Bob}, e \textbf{Eve} (l'avversaria, da \emph{eavesdropper}) sta in mezzo al canale.

Due proprietà che vogliamo:

\begin{itemize}
\tightlist
\item
  \textbf{Confidenzialità}: Eve intercetta \(m\) ma non deve capirne nulla.
\item
  \textbf{Integrità}: Eve non deve poter modificare il messaggio. Se Bob riceve \(m' \ne m\) deve accorgersene.
\end{itemize}

\begin{nota}

In questa lezione ci occupiamo solo di \textbf{confidenzialità}.

\end{nota}

\section{Sicurezza incondizionata (Shannon)}

\begin{inbreve}

La sicurezza è ``incondizionata'' quando vale anche contro un avversario con potenza di calcolo infinita. Non si basa su problemi difficili da calcolare, ma sul fatto che l'informazione per rompere lo schema proprio non c'è.

\end{inbreve}

È la trattazione dovuta a \textbf{Claude Shannon}. ``Incondizionata'' vuol dire che la sicurezza vale anche contro un avversario con \textbf{potenza di calcolo illimitata}.

\subsection{Schema di cifratura}

Alice e Bob condividono una chiave segreta \(k\). Alice cifra \(m\) ottenendo il cifrato \(c\), lo manda a Bob, e Bob decifra con la stessa chiave.

\begin{verbatim}
         k                          k
         ↓                          ↓
Alice → [Enc] ──── c ────→ [Dec] → Bob
  m
\end{verbatim}

Notazione:

\begin{itemize}
\tightlist
\item
  \(\mathcal{K}\) = spazio delle chiavi, \(\mathcal{M}\) = spazio dei messaggi, \(\mathcal{C}\) = spazio dei cifrati
\item
  Uno schema è una coppia \(\Pi = (\text{Enc}, \text{Dec})\) con
\end{itemize}

\[
\text{Enc}: \mathcal{K} \times \mathcal{M} \to \mathcal{C} \qquad \text{Dec}: \mathcal{K} \times \mathcal{C} \to \mathcal{M}
\]

\textbf{Ipotesi}: la chiave \(k \in \mathcal{K}\) è scelta in modo \textbf{uniformemente casuale} ed è \textbf{sconosciuta a Eve}.

\subsection{Correttezza}

Decifrare quello che si è cifrato deve restituire il messaggio originale:

\[
\forall k \in \mathcal{K},\ \forall m \in \mathcal{M}: \quad \text{Dec}(k, \text{Enc}(k, m)) = m
\]

\subsection{Definizione: segretezza perfetta}

\begin{inbreve}

Uno schema è perfettamente segreto se vedere il cifrato non cambia nulla di quello che Eve sa sul messaggio: la sua stima prima e dopo aver visto il cifrato è la stessa.

\end{inbreve}

Sia \(M\) una \textbf{qualunque} distribuzione sui messaggi e \(K\) la distribuzione uniforme sulle chiavi. Allora anche il cifrato \(C = \text{Enc}(K, M)\) è una variabile aleatoria.

\(\Pi\) è \textbf{perfettamente segreto} se per ogni distribuzione \(M\), per ogni \(m \in \mathcal{M}\) e per ogni \(c \in \mathcal{C}\) (con \(\Pr[C=c] > 0\)):

\[
\Pr[M = m] = \Pr[M = m \mid C = c]
\]

\textbf{Intuizione}: \(\Pr[M = m]\) è quello che Eve sa sul messaggio \emph{prima} di vedere il cifrato (es. ``Alice scrive `sì' nel 70\% dei casi''). \(\Pr[M = m \mid C = c]\) è quello che sa \emph{dopo} averlo visto. Se sono uguali, vedere il cifrato non le ha insegnato nulla.

\begin{nota}

\textbf{Perché ``per ogni distribuzione} \(M\)\textbf{``?} Perché non sappiamo cosa Eve sappia già dei messaggi. La definizione deve funzionare qualunque sia la sua conoscenza a priori.

\end{nota}

\section{Teorema: tre formulazioni equivalenti}

\begin{inbreve}

La segretezza perfetta si può dire in tre modi equivalenti: (i) a priori = a posteriori, (ii) messaggio e cifrato sono indipendenti, (iii) ogni messaggio produce i cifrati con la stessa distribuzione. Per dimostrare che uno schema è perfettamente segreto di solito si usa la (iii).

\end{inbreve}

Le seguenti sono equivalenti:

\begin{enumerate}
\def\labelenumi{\arabic{enumi}.}
\tightlist
\item
  \textbf{(i)} \(\Pi\) è perfettamente segreto
\item
  \textbf{(ii)} \(M\) e \(C\) sono \textbf{indipendenti} (in termini di teoria dell'informazione: la mutua informazione è \(I(M; C) = 0\))
\item
  \textbf{(iii)} \(\forall m, m' \in \mathcal{M},\ \forall c \in \mathcal{C}\):
\end{enumerate}

\[
\Pr[\text{Enc}(K, m) = c] = \Pr[\text{Enc}(K, m') = c]
\]

\textbf{Intuizione di (iii)}: la distribuzione del cifrato è \textbf{la stessa qualunque sia il messaggio}. Vedendo \(c\) non puoi capire se veniva da \(m\) o da \(m'\). Nota che qui la probabilità è solo sulla scelta della chiave: il messaggio è fissato.

\begin{enumerate}
\def\labelenumi{(\roman{enumi})}
\setcounter{enumi}{2}
\tightlist
\item
  è la formulazione più comoda da verificare in pratica, ed è quella che si usa per dimostrare che uno schema è perfettamente segreto.
\end{enumerate}

\section{One-Time Pad (OTP)}

\textbf{Definizione}: il One-Time Pad è un cifrario a chiave simmetrica che cifra facendo lo \textbf{XOR bit a bit} tra messaggio e chiave. La chiave è scelta \textbf{uniformemente}, è \textbf{lunga quanto il messaggio} e va usata \textbf{una sola volta}.

\(\mathcal{K} = \mathcal{M} = \mathcal{C} = \{0,1\}^\ell\) (stringhe di \(\ell\) bit)

\[
\text{Enc}(k, m) = k \oplus m \qquad \text{Dec}(k, c) = k \oplus c
\]

\textbf{Correttezza}: \(\text{Dec}(k, \text{Enc}(k, m)) = k \oplus (k \oplus m) = (k \oplus k) \oplus m = 0^\ell \oplus m = m\) ✓

\begin{esempio}

\textbf{Esempio} con \(\ell = 4\): \(m = 1010\), \(k = 0110\)

Cifratura: \(c = 1010 \oplus 0110 = 1100\)

Decifratura: \(1100 \oplus 0110 = 1010 = m\) ✓

\end{esempio}

\subsection{Corollario: l'OTP è perfettamente segreto}

\begin{inbreve}

Per ogni messaggio e ogni cifrato c'è esattamente una chiave che porta l'uno nell'altro. Siccome la chiave è casuale, visto il cifrato ogni messaggio resta ugualmente possibile.

\end{inbreve}

\textbf{Dimostrazione} (usando la (iii)). Fissiamo un qualunque \(m\) e un qualunque \(c\):

\[
\Pr[\text{Enc}(K, m) = c] = \Pr[K \oplus m = c] = \Pr[K = c \oplus m] = 2^{-\ell}
\]

\begin{itemize}
\tightlist
\item
  Il secondo passaggio: ``\(K \oplus m = c\)'' equivale a ``\(K = c \oplus m\)'' (fai XOR con \(m\) da entrambi i lati).
\end{itemize}

È come risolvere un'equazione normale. Se hai x + 3 = 5, per isolare x sottrai 3 da entrambi i lati. Con lo XOR si fa la stessa cosa, ma l'operazione ``inversa'' dello XOR è lo XOR stesso: per togliere m, fai di nuovo ⊕ m.

Passo per passo, partendo da

K ⊕ m = c

fai ⊕ m da entrambi i lati:

K ⊕ m ⊕ m = c ⊕ m

A sinistra m ⊕ m = 0 (una cosa XOR se stessa si annulla), quindi resta

K ⊕ 0 = c ⊕ m

e siccome K ⊕ 0 = K:

K = c ⊕ m

\begin{itemize}
\tightlist
\item
  L'ultimo: \(c \oplus m\) è \textbf{una} stringa specifica, e \(K\) è uniforme su \(2^\ell\) stringhe, quindi la probabilità di beccarla è \(2^{-\ell}\).
\end{itemize}

Lo stesso conto per un qualunque altro \(m'\) dà \(\Pr[\text{Enc}(K, m') = c] = 2^{-\ell}\). Le due probabilità sono uguali, quindi vale (iii) e per il teorema l'OTP è perfettamente segreto. ∎

\textbf{Intuizione}: per ogni coppia \((m, c)\) esiste \textbf{esattamente una} chiave che porta \(m\) in \(c\). Quindi, visto \(c\), ogni messaggio è ugualmente plausibile: basta ``supporre'' la chiave giusta.

\subsection{Limiti dell'OTP}

La segretezza perfetta è possibile, ma \textbf{a caro prezzo}:

\begin{enumerate}
\def\labelenumi{\arabic{enumi}.}
\tightlist
\item
  \textbf{La chiave è lunga quanto il messaggio.} Per mandare 1 GB serve 1 GB di chiave condivisa in anticipo.
\item
  \textbf{La chiave si può usare una sola volta} (da cui \emph{one-time}). Se cifri due messaggi con la stessa chiave:
\end{enumerate}

\[
c_1 = k \oplus m_1, \quad c_2 = k \oplus m_2 \quad\Rightarrow\quad c_1 \oplus c_2 = m_1 \oplus m_2
\]

perché \(k \oplus k = 0^\ell\). La chiave sparisce ed Eve ottiene lo XOR dei due messaggi.

\begin{esempio}

\textbf{Esempio}: \(m_1 \oplus m_2\) vale 0 esattamente nelle posizioni in cui i due messaggi coincidono. Eve scopre quindi dove sono uguali e dove no. E se conosce uno dei due messaggi, ricava l'altro: \(m_2 = (c_1 \oplus c_2) \oplus m_1\).

\end{esempio}

\textbf{Perché la chiave sparisce}: lo XOR di un bit con se stesso fa sempre 0, quindi facendo \(c_1 \oplus c_2\) le due \(k\) si annullano:

\[
c_1 \oplus c_2 = (m_1 \oplus k) \oplus (m_2 \oplus k) = m_1 \oplus m_2 \oplus (k \oplus k) = m_1 \oplus m_2
\]

\begin{esempio}

\textbf{Esempio numerico}: chiave \(k = 0110\) usata due volte.

\begin{itemize}
\item
  \(m_1 = 1010\), quindi \(c_1 = 1010 \oplus 0110 = 1100\)
\item
  \(m_2 = 1001\), quindi \(c_2 = 1001 \oplus 0110 = 1111\)
\end{itemize}

Eve vede solo \(c_1\) e \(c_2\) e calcola \(c_1 \oplus c_2 = 1100 \oplus 1111 = 0011\), che è proprio \(m_1 \oplus m_2 = 1010 \oplus 1001 = 0011\), \textbf{senza conoscere la chiave}.

\end{esempio}

\begin{nota}

\textbf{Cosa ci legge}: 0011 vuol dire che i due messaggi hanno i primi due bit uguali e gli ultimi due diversi.

\end{nota}

\begin{nota}

\textbf{Se conosce} \(m_1 = 1010\): ricava \(m_2 = 0011 \oplus 1010 = 1001\) e anche la chiave \(k = c_1 \oplus m_1 = 1100 \oplus 1010 = 0110\). Nella realtà conoscere un pezzo di messaggio è facile: le email iniziano con ``Ciao'', i file hanno intestazioni standard.

\end{nota}

\textbf{Il confronto}: con \textbf{un solo} cifrato ogni messaggio è compatibile con \(c\) (basta cambiare chiave), quindi Eve non impara nulla. Con \textbf{due} cifrati e la stessa chiave, combinandoli elimina la chiave e ottiene informazioni vere sui messaggi.

\section{Dimostrazione del teorema}

Schema: (i) ⇒ (ii) ⇒ (iii) ⇒ (i). Così basta dimostrare tre implicazioni per avere tutte le equivalenze.

\subsection{(i) ⇒ (ii)}

Per la segretezza perfetta e la definizione di probabilità condizionata:

\[
\Pr[M = m] = \Pr[M = m \mid C = c] = \frac{\Pr[M = m \wedge C = c]}{\Pr[C = c]}
\]

Moltiplicando per \(\Pr[C = c]\):

\[
\Pr[M = m \wedge C = c] = \Pr[M = m] \cdot \Pr[C = c]
\]

che è proprio la definizione di indipendenza tra \(M\) e \(C\).

\subsection{(ii) ⇒ (iii)}

\[
\Pr[\text{Enc}(K, m) = c] = \Pr[\text{Enc}(K, M) = c \mid M = m] = \Pr[C = c \mid M = m] = \Pr[C = c]
\]

\begin{itemize}
\tightlist
\item
  \textbf{1° passaggio}: condizionare su \(M = m\) equivale a fissare il messaggio a \(m\). Funziona perché la chiave \(K\) è scelta indipendentemente dal messaggio, quindi sapere \(M = m\) non cambia la distribuzione di \(K\).
\item
  \textbf{2° passaggio}: \(\text{Enc}(K, M)\) è per definizione \(C\).
\item
  \textbf{3° passaggio}: per (ii) \(M\) e \(C\) sono indipendenti, quindi il condizionamento sparisce.
\end{itemize}

Con gli stessi passaggi, \(\Pr[\text{Enc}(K, m') = c] = \Pr[C = c]\). Entrambe valgono \(\Pr[C = c]\), quindi sono uguali. ✓

\subsection{(iii) ⇒ (i)}

\textbf{Passo 1}: mostriamo che \(\Pr[C = c \mid M = m] = \Pr[C = c]\).

Partiamo da \(\Pr[C = c]\) e usiamo la \textbf{probabilità totale} spezzando su tutti i possibili messaggi \(m'\):

\[
\begin{aligned}
\Pr[C = c] &= \sum_{m'} \Pr[C = c \wedge M = m'] \\
&= \sum_{m'} \Pr[C = c \mid M = m'] \cdot \Pr[M = m'] \\
&= \sum_{m'} \Pr[\text{Enc}(K, M) = c \mid M = m'] \cdot \Pr[M = m'] \\
&= \sum_{m'} \Pr[\text{Enc}(K, m') = c] \cdot \Pr[M = m'] \\
&= \sum_{m'} \Pr[\text{Enc}(K, m) = c] \cdot \Pr[M = m'] \qquad \text{(qui usiamo (iii))} \\
&= \Pr[\text{Enc}(K, m) = c] \cdot \underbrace{\sum_{m'} \Pr[M = m']}_{=\,1} \\
&= \Pr[\text{Enc}(K, M) = c \mid M = m] = \Pr[C = c \mid M = m]
\end{aligned}
\]

\textbf{Il trucco}: grazie a (iii), \(\Pr[\text{Enc}(K, m') = c]\) è lo stesso numero per ogni \(m'\), quindi lo sostituiamo con quello di un \(m\) fissato e lo portiamo fuori dalla somma. Quel che resta è la somma delle probabilità di tutti i messaggi, che fa 1.

\textbf{Passo 2}: applichiamo Bayes. Scriviamo \(\Pr[M = m \wedge C = c]\) in due modi:

\[
\Pr[M = m \mid C = c] \cdot \Pr[C = c] = \Pr[M = m \wedge C = c] = \Pr[C = c \mid M = m] \cdot \Pr[M = m]
\]

Quindi

\[
\Pr[M = m] = \frac{\Pr[M = m \mid C = c] \cdot \Pr[C = c]}{\Pr[C = c \mid M = m]}
\]

Per il passo 1 il denominatore è uguale a \(\Pr[C = c]\), che si semplifica con quello al numeratore:

\[
\Pr[M = m] = \Pr[M = m \mid C = c]
\]

che è la segretezza perfetta. ∎

\section{Riepilogo}

\begin{itemize}
\tightlist
\item
  Sicurezza dimostrabile = definizione precisa + dimostrazione
\item
  Confidenzialità (Eve non capisce) vs integrità (Eve non modifica)
\item
  Segretezza perfetta: vedere il cifrato non cambia la conoscenza sul messaggio, \(\Pr[M = m] = \Pr[M = m \mid C = c]\)
\item
  Equivale a \(M\) e \(C\) indipendenti, ed equivale a ``ogni messaggio produce la stessa distribuzione di cifrati''
\item
  L'OTP (\(c = k \oplus m\)) è perfettamente segreto
\item
  Prezzo: chiave lunga quanto il messaggio e usabile una volta sola
\end{itemize}

% Generato da tools/notion2latex.py a partire dalla pagina Notion "lezione-2".

\chapter{Lezione 2 -- Teorema di Shannon e MAC}

\section{Teorema di Shannon: il prezzo della segretezza perfetta}

Nella lezione 1 abbiamo visto che l'OTP è perfettamente segreto ma ha limiti pesanti. Shannon dimostra che non è colpa dell'OTP: \textbf{ogni} schema perfettamente segreto deve avere gli stessi limiti.

\subsection{Teorema}

\begin{inbreve}

Per avere segretezza perfetta servono almeno tante chiavi quanti messaggi, cioè una chiave lunga almeno quanto il messaggio. Vale per qualunque cifrario, non solo per l'OTP.

\end{inbreve}

Per ogni SKE (schema di cifratura a chiave simmetrica) con segretezza perfetta:

\[
|\mathcal{K}| \ge |\mathcal{M}|
\]

Cioè servono almeno tante chiavi quanti sono i messaggi. Con chiavi e messaggi fatti di bit: la chiave deve essere lunga \textbf{almeno quanto il messaggio}.

\subsection{Dimostrazione (per assurdo)}

\begin{inbreve}

Se le chiavi sono meno dei messaggi, Eve prova a decifrare il cifrato con tutte le chiavi e ottiene meno candidati dei messaggi possibili. Qualche messaggio viene escluso, quindi Eve ha imparato qualcosa.

\end{inbreve}

Sia \(\Pi\) un qualunque SKE perfettamente segreto. Fissiamo la \textbf{distribuzione uniforme} su \(\mathcal{M}\) (la definizione deve valere per ogni distribuzione, quindi possiamo sceglierne una comoda). Prendiamo un qualunque \(c \in \mathcal{C}\) con \(\Pr[C = c] > 0\).

Consideriamo l'insieme di tutti i messaggi che \(c\) potrebbe nascondere:

\[
\mathcal{M}' = \{\text{Dec}(k, c) : k \in \mathcal{K}\}
\]

Supponiamo per assurdo che \(|\mathcal{K}| < |\mathcal{M}|\).

\begin{itemize}
\tightlist
\item
  \(|\mathcal{M}'| \le |\mathcal{K}|\): ogni chiave produce un solo messaggio (Dec è deterministica), quindi i messaggi diversi in \(\mathcal{M}'\) sono al massimo quante le chiavi.
\item
  Quindi \(|\mathcal{M}'| \le |\mathcal{K}| < |\mathcal{M}|\), cioè \(|\mathcal{M}'| < |\mathcal{M}|\).
\end{itemize}

Allora esiste un messaggio \(m \in \mathcal{M} \setminus \mathcal{M}'\): un messaggio che \(c\) non può nascondere con nessuna chiave.

\begin{nota}

\textbf{Attenzione alla notazione}: nelle formule compaiono tre simboli simili, con significati diversi.

\begin{itemize}
\tightlist
\item
  \(\mathcal{M}\) (corsiva) è l'\textbf{insieme} di tutti i messaggi possibili; \(\mathcal{M}'\) è il suo sottoinsieme dei candidati per \(c\).
\item
  \(M\) (dritta), dentro \(\Pr[M = m]\), \textbf{non è un insieme}: è ``il messaggio che Alice manda davvero'', che Eve non conosce (una variabile aleatoria, come ``il risultato del dado'').
\item
  \(m\) (minuscola) è \textbf{un singolo messaggio}: quello escluso, che sta in \(\mathcal{M}\) ma non in \(\mathcal{M}'\).
\end{itemize}

Quindi \(\Pr[M = m]\) si legge ``la probabilità che Alice abbia mandato proprio \(m\)'', e riguarda solo quel messaggio, non tutto l'insieme.

\end{nota}

\begin{verbatim}
┌───────────────────────────┐
│ M            × m          │
│        ┌──────────┐       │
│        │    M'    │       │
│        └──────────┘       │
└───────────────────────────┘
\end{verbatim}

Ora confrontiamo:

\begin{itemize}
\tightlist
\item
  \textbf{A priori}: \(\Pr[M = m] = 1/|\mathcal{M}| > 0\) (distribuzione uniforme)
\item
  \textbf{A posteriori}: \(\Pr[M = m \mid C = c] = 0\). Infatti, se \(c\) fosse la cifratura di \(m\) con qualche chiave \(k\), per correttezza \(\text{Dec}(k, c) = m\), e quindi \(m\) starebbe in \(\mathcal{M}'\). Quindi \textbf{nessuna chiave} trasforma \(m\) in \(c\), e con la formula della probabilità condizionata:
\end{itemize}

\[
\Pr[M = m \mid C = c] = \frac{\Pr[M = m \wedge C = c]}{\Pr[C = c]} = \frac{0}{\Pr[C = c]} = 0
\]

Il numeratore è 0 perché ``Alice manda \(m\) \textbf{e} il cifrato è \(c\)'' non succede mai; il denominatore è positivo perché abbiamo scelto un \(c\) con \(\Pr[C = c] > 0\).

A priori e a posteriori sono diversi: contraddizione con la segretezza perfetta. ∎

\textbf{Intuizione}: con meno chiavi che messaggi, Eve può prendere \(c\) e decifrarlo con tutte le chiavi possibili. I messaggi che non escono mai sono esclusi: vedere \(c\) le ha insegnato qualcosa.

\subsection{Takeaway}

La segretezza perfetta costa troppo, quindi \textbf{faremo qualcosa di diverso}.

\textbf{Obiettivo}: Alice e Bob condividono \textbf{una} chiave di lunghezza \textbf{indipendente} dalla lunghezza del messaggio. \emph{(Per riuscirci bisognerà rinunciare alla sicurezza contro avversari con potenza di calcolo illimitata.)}

\section{Esercizio: segretezza perfetta come gioco}

\begin{inbreve}

Un altro modo di definire la segretezza perfetta. Si sfida Eve a capire quale tra due messaggi è stato cifrato: lo schema è sicuro se nessuna strategia di Eve si comporta diversamente a seconda del messaggio cifrato.

\end{inbreve}

Definizione alternativa di segretezza perfetta tramite un \textbf{gioco} (esperimento) tra un \textbf{avversario} \(A\) e un \textbf{challenger} \(C\).

\(\text{GAME}_A(b, m_0, m_1)\) con \(b \in \{0,1\}\) e \(m_0, m_1 \in \mathcal{M}\):

\begin{enumerate}
\def\labelenumi{\arabic{enumi}.}
\tightlist
\item
  \(C\) sceglie \(k \leftarrow \mathcal{K}\) uniformemente
\item
  \(C\) manda ad \(A\) il cifrato \(c = \text{Enc}(k, m_b)\)
\item
  \(A\) risponde con \(b' \in \{0,1\}\), il suo tentativo di indovinare quale messaggio è stato cifrato
\end{enumerate}

\begin{verbatim}
Avversario A                         Challenger C
                                     k ← K
      ←──────── c = Enc(k, m_b) ────────
      ──────── b' ∈ {0,1} ─────────→
\end{verbatim}

\textbf{Definizione}: \(\Pi\) è perfettamente segreto se \(\forall A\), \(\forall m_0, m_1 \in \mathcal{M}\):

\[
\Big| \Pr[b' = 1 : \text{GAME}_A(0, m_0, m_1)] - \Pr[b' = 1 : \text{GAME}_A(1, m_0, m_1)] \Big| = 0
\]

Notazione: \(\Pr[b' = 1 : \text{GAME}]\) è la probabilità che \(A\) risponda 1 giocando quel gioco.

\textbf{Formulazione alternativa}: \(\text{GAME}_A(0, m_0, m_1) \equiv \text{GAME}_A(1, m_0, m_1)\), cioè i due giochi sono \textbf{identicamente distribuiti}.

\textbf{Intuizione}: se la risposta di \(A\) dipendesse da \(b\), vorrebbe dire che \(A\) riesce a distinguere la cifratura di \(m_0\) da quella di \(m_1\). La definizione chiede che nessun avversario ci riesca.

\subsection{Spiegazione semplice}

\textbf{L'idea}: mettiamo alla prova Eve. Le facciamo vedere un cifrato e le chiediamo \textbf{quale di due messaggi} c'è dentro. Se il cifrario è sicuro, non deve riuscire a capirlo.

\textbf{Le regole}: si scelgono due messaggi, per esempio \(m_0\) = ``SÌ'' e \(m_1\) = ``NO''. Poi:

\begin{enumerate}
\def\labelenumi{\arabic{enumi}.}
\tightlist
\item
  \(C\) (l'arbitro) sceglie una chiave a caso.
\item
  \(C\) cifra \textbf{uno dei due} messaggi e dà il cifrato \(c\) ad \(A\) (Eve).
\item
  \(A\) guarda \(c\) e risponde \textbf{0} (``secondo me era \(m_0\)'') oppure \textbf{1} (``secondo me era \(m_1\)'').
\end{enumerate}

\textbf{GAME(0) e GAME(1)}: la \(b\) dice quale messaggio viene cifrato. In \textbf{GAME(0)} \(C\) cifra sempre \(m_0\), in \textbf{GAME(1)} cifra sempre \(m_1\). \(A\) non sa in quale versione si trova.

\textbf{La condizione di sicurezza}: il cifrario è sicuro se \(A\) si comporta \textbf{allo stesso modo} in entrambe le versioni, cioè risponde ``1'' con la stessa probabilità sia in GAME(0) sia in GAME(1). Se le sue risposte cambiano a seconda del messaggio cifrato, vuol dire che dal cifrato riesce a capire qualcosa.

\begin{esempio}

\textbf{Esempio: cifrario rotto.} Chiave di 1 bit, messaggi di 2 bit: se \(k = 1\) si invertono entrambi i bit, se \(k = 0\) restano uguali. Prendiamo \(m_0 = 00\) e \(m_1 = 01\).

\begin{itemize}
\item
  In \textbf{GAME(0)} si cifra 00, quindi \(c\) è 00 oppure 11: bit sempre \textbf{uguali}.
\item
  In \textbf{GAME(1)} si cifra 01, quindi \(c\) è 01 oppure 10: bit sempre \textbf{diversi}.
\end{itemize}

\(A\) usa la strategia ``se i bit di \(c\) sono uguali rispondo 0, se sono diversi rispondo 1''. In GAME(0) risponde 1 con probabilità \textbf{0}, in GAME(1) con probabilità \textbf{1}. La differenza è 1: \(A\) vince sempre, quindi il cifrario \textbf{non} è sicuro.

\end{esempio}

\begin{esempio}

\textbf{Esempio: OTP.} Chiave di 2 bit, sempre \(m_0 = 00\) e \(m_1 = 01\).

\begin{itemize}
\item
  In \textbf{GAME(0)} \(c = 00 \oplus k\): uno qualsiasi tra 00, 01, 10, 11, ognuno con probabilità 1/4.
\item
  In \textbf{GAME(1)} \(c = 01 \oplus k\): di nuovo uno qualsiasi tra 00, 01, 10, 11, ognuno con probabilità 1/4.
\end{itemize}

\(A\) vede cifrati distribuiti \textbf{esattamente allo stesso modo} nei due giochi, quindi qualunque strategia usi risponde 1 con la stessa probabilità in entrambi. La differenza è 0: l'OTP è sicuro.

\end{esempio}

\textbf{In una frase}: il gioco chiede ``Eve, guardando il cifrato, riesce a distinguere se dentro c'era \(m_0\) o \(m_1\)?''. Se per ogni coppia di messaggi e per ogni strategia la risposta è no, il cifrario è perfettamente segreto. È la stessa idea di ``a priori = a posteriori'', detta in un altro modo.

\textbf{Esercizio}: dimostrare che questa definizione è equivalente alla segretezza perfetta.

\begin{approfondimento}{Soluzione proposta (non negli appunti del prof)}

\textbf{PS ⇒ gioco.} Per la (iii) della lezione 1, \(\text{Enc}(K, m_0)\) e \(\text{Enc}(K, m_1)\) hanno la stessa distribuzione. \(A\) vede solo \(c\), quindi nei due giochi riceve un input con la stessa distribuzione e la sua risposta \(b'\) ha la stessa distribuzione. La differenza è 0.

\textbf{Gioco ⇒ PS} (per contrapposizione). Se \(\Pi\) non è perfettamente segreto, per la (iii) esistono \(m_0, m_1, c\) con \(\Pr[\text{Enc}(K, m_0) = c] \ne \Pr[\text{Enc}(K, m_1) = c]\). Costruiamo l'avversario \(A\): ``se ricevo \(c\) rispondo 1, altrimenti 0''. Allora \(\Pr[b' = 1 : \text{GAME}_A(0)] = \Pr[\text{Enc}(K, m_0) = c]\) e \(\Pr[b' = 1 : \text{GAME}_A(1)] = \Pr[\text{Enc}(K, m_1) = c]\), che sono diverse: la differenza non è 0 e la definizione col gioco è violata.

\end{approfondimento}

\section{Message Authentication Codes (MAC)}

\begin{inbreve}

Un MAC serve a \textbf{autenticare} un messaggio, cioè a garantire che venga davvero da chi conosce la chiave e che non sia stato modificato. Alice attacca al messaggio un \textbf{tag} che solo chi ha la chiave sa calcolare, e Bob lo ricalcola per verificarlo. Non nasconde il messaggio.

\end{inbreve}

Nella lezione 1 ci siamo occupati di \textbf{confidenzialità}. Ora passiamo all'\textbf{integrità}: Eve non deve poter modificare il messaggio senza che Bob se ne accorga.

\subsection{Schema}

Alice e Bob condividono una chiave \(k \in \mathcal{K}\).

\begin{verbatim}
Alice (k)                                       Bob (k)
m → [Tag] → τ                                   m → [Tag] → τ*
                ────── (m, τ) ──────→                  τ* = τ ?
\end{verbatim}

\begin{itemize}
\tightlist
\item
  Alice calcola il \textbf{tag} \(\tau = \text{Tag}(k, m)\) e manda la coppia \((m, \tau)\), \textbf{in chiaro}: il MAC non nasconde il messaggio, lo autentica.
\item
  Bob ricalcola \(\text{Tag}(k, m)\) e controlla che coincida con \(\tau\).
\end{itemize}

\[
\text{Tag}: \mathcal{K} \times \mathcal{M} \to \mathcal{T}
\]

Tag è \textbf{deterministica}: stessa chiave e stesso messaggio danno sempre lo stesso tag. È quello che permette a Bob di ricalcolarlo e confrontarlo.

\subsection{Sicurezza: cosa vuole fare Eve}

\textbf{Obiettivo di Eve}: produrre una coppia \((m', \tau')\), \textbf{senza conoscere} \(k\), tale che \(\tau' = \text{Tag}(k, m')\). Si chiama \emph{forgery} (falsificazione).

Rendiamo Eve più forte: conosce già una coppia \textbf{valida} \((m, \tau)\), cioè con \(\text{Tag}(k, m) = \tau\) (ad esempio perché l'ha intercettata). Deve comunque essere difficile falsificare un \(\tau'\) valido per un \(m' \ne m\).

\begin{nota}

Perché \(m' \ne m\)? Rimandare la stessa coppia \((m, \tau)\) è sempre possibile, quindi non conta come falsificazione.

\end{nota}

\subsection{Definizione: sicurezza statistica (one-time)}

\begin{inbreve}

Un MAC è ε-sicuro se Eve, anche dopo aver visto un messaggio con il suo tag, riesce a indovinare il tag di un altro messaggio al massimo con probabilità \(\varepsilon\), piccolissima.

\end{inbreve}

Sia Tag un MAC su \(\mathcal{K}, \mathcal{M}, \mathcal{T}\). Diciamo che è \textbf{ε-statisticamente sicuro} (\emph{one-time}) se \(\forall m, m' \in \mathcal{M}\) con \(m \ne m'\) e \(\forall \tau, \tau' \in \mathcal{T}\):

\[
\Pr[\text{Tag}(K, m') = \tau' \mid \text{Tag}(K, m) = \tau] \le \varepsilon
\]

con \(K\) uniforme su \(\mathcal{K}\).

\textbf{Intuizione}: anche sapendo il tag di \(m\), indovinare il tag di un altro messaggio riesce al massimo con probabilità \(\varepsilon\).

\begin{itemize}
\tightlist
\item
  \textbf{Statistica}: vale contro avversari con potenza di calcolo illimitata, come la segretezza perfetta di Shannon.
\item
  \textbf{One-time}: Eve vede una sola coppia valida, quindi la chiave va usata per autenticare un solo messaggio.
\end{itemize}

Valori tipici di \(\varepsilon\): \(2^{-80}\), \(2^{-128}\).

\subsection{Esercizi}

\textbf{Esercizio 1.} La ε-sicurezza è impossibile per \(\varepsilon = 0\).

\begin{approfondimento}{Soluzione proposta (non negli appunti del prof)}

Fissiamo \(m \ne m'\) e un \(\tau\) con \(\Pr[\text{Tag}(K, m) = \tau] > 0\). Il tag di \(m'\) sarà pur qualcosa, quindi sommando su tutti i \(\tau' \in \mathcal{T}\) le probabilità condizionate \(\Pr[\text{Tag}(K, m') = \tau' \mid \text{Tag}(K, m) = \tau]\) si ottiene 1. Sono \(|\mathcal{T}|\) termini, quindi almeno uno vale \(\ge 1/|\mathcal{T}| > 0\). Ne segue \(\varepsilon \ge 1/|\mathcal{T}|\).

In pratica: Eve può sempre tirare a indovinare un tag a caso.

\end{approfondimento}

\textbf{Esercizio 2.} L'OTP usato come MAC è insicuro: \(\text{Tag}(k, m) = \tau = m \oplus k\).

\begin{approfondimento}{Soluzione proposta (non negli appunti del prof)}

Eve conosce \((m, \tau)\) e ricava la chiave: \(k = m \oplus \tau\). Poi per qualunque \(m' \ne m\) calcola \(\tau' = m' \oplus k\). Riesce con probabilità 1, quindi lo schema non è ε-sicuro per nessun \(\varepsilon < 1\).

Esempio: \(m = 1010\), \(\tau = 0011\), quindi \(k = 1001\). Per \(m' = 1111\): \(\tau' = 1111 \oplus 1001 = 0110\).

Morale: l'OTP è perfetto per la confidenzialità ma inutile per l'integrità.

\end{approfondimento}

\textbf{Esercizio 3.} Consideriamo un gioco come prima:

\begin{enumerate}
\def\labelenumi{\arabic{enumi}.}
\tightlist
\item
  \(A\) manda un messaggio \(m \in \mathcal{M}\) al challenger
\item
  Il challenger sceglie \(k \leftarrow \mathcal{K}\) e risponde con \(\tau = \text{Tag}(k, m)\)
\item
  \(A\) produce una coppia \((m', \tau')\)
\end{enumerate}

\begin{verbatim}
Avversario A                         Challenger C
      ──────── m ∈ M ──────────→
                                     k ← K
      ←──────── τ = Tag(k, m) ───────
      ──────── (m', τ') ─────────→
\end{verbatim}

\(A\) \textbf{vince} (WIN) se: (i) \(m' \ne m\) e (ii) \(\text{Tag}(k, m') = \tau'\).

Richiesta: \(\forall A\), \(\Pr[\text{WIN}] \le \varepsilon\).

\begin{approfondimento}{Soluzione proposta (non negli appunti del prof)}

\textbf{Definizione ⇒ gioco.} Per un \(A\) deterministico, \(m\) è fissato e la coppia \((m', \tau')\) dipende solo dal \(\tau\) ricevuto. Spezzando sui possibili \(\tau\) (probabilità totale): \(\Pr[\text{WIN}] = \sum_\tau \Pr[\text{Tag}(K, m) = \tau] \cdot \Pr[\text{Tag}(K, m') = \tau' \mid \text{Tag}(K, m) = \tau]\). Ogni probabilità condizionata è al più \(\varepsilon\) e le altre sommano a 1, quindi \(\Pr[\text{WIN}] \le \varepsilon\). Se \(A\) è probabilistico, si fa la media sulle sue scelte casuali e il limite resta.

\textbf{Il viceversa in generale no}: la definizione chiede il limite per \textbf{ogni} \(\tau\), il gioco solo \textbf{in media} sui \(\tau\). Un tag rarissimo che rende facile la falsificazione viola la definizione ma pesa pochissimo nel gioco.

\end{approfondimento}

\section{Costruire MAC: funzioni hash pairwise independent}

\begin{inbreve}

Si sceglie a caso una funzione da una famiglia (cioè si sceglie una chiave a caso). Ogni input dà \textbf{il suo} output. La famiglia è \textbf{pairwise independent} se, \textbf{conoscendo l'output su un input, non si sa nulla dell'output su un altro input}: può essere qualunque valore, compreso lo stesso, ognuno con la stessa probabilità. È come lanciare due dadi separati: a volte escono uguali, ma il primo non dice niente sul secondo. Nel MAC l'output è il tag: Eve vede il tag di un messaggio e non impara nulla sul tag di un altro.

\end{inbreve}

\begin{esempio}

\textbf{Esempio}: \(h(m) = a \cdot m + b \bmod 3\). Se sai \(h(0) = 2\), allora \(b = 2\), ma \(a\) resta ignoto: \(h(1) = a + 2\) vale 2, 0 o 1 a seconda di \(a\), ognuno con probabilità 1/3. Con \(h(m) = m + k \bmod 3\) invece, da \(h(0) = 2\) ricavi \(k = 2\) e quindi \(h(1) = 0\) con certezza: \textbf{non} è pairwise independent.

\end{esempio}

Costruiremo i MAC usando un particolare tipo di \textbf{funzioni hash}.

\subsection{Definizione}

Una famiglia \(\mathcal{H} = \{h_k : \mathcal{M} \to \mathcal{T}\}_{k \in \mathcal{K}}\) è \textbf{pairwise independent} se \(\forall m, m' \in \mathcal{M}\) con \(m \ne m'\), la coppia

\[
(h(K, m),\ h(K, m'))
\]

con \(K\) uniforme su \(\mathcal{K}\), è \textbf{uniforme} su \(\mathcal{T}^2 = \mathcal{T} \times \mathcal{T}\).

Notazione: \(h(k, m) = h_k(m)\).

\textbf{Intuizione}: i valori della funzione su due input diversi si comportano come due valori scelti a caso e indipendenti. Conoscere \(h_K(m)\) non dice nulla su \(h_K(m')\).

Questo suggerisce di usare:

\[
\text{Tag}(k, m) = h(k, m) = h_k(m)
\]

\begin{approfondimento}{Perché funziona (anticipazione, non negli appunti del prof)}

Se la coppia è uniforme su \(\mathcal{T}^2\), ogni coppia \((\tau, \tau')\) ha probabilità \(1/|\mathcal{T}|^2\), e da sola \(h(K, m) = \tau\) ha probabilità \(1/|\mathcal{T}|\). Quindi:

\(\Pr[h(K, m') = \tau' \mid h(K, m) = \tau] = \dfrac{1/|\mathcal{T}|^2}{1/|\mathcal{T}|} = \dfrac{1}{|\mathcal{T}|}\)

Il MAC è ε-sicuro con \(\varepsilon = 1/|\mathcal{T}|\), il minimo possibile per l'esercizio 1. Con tag di 128 bit: \(\varepsilon = 2^{-128}\).

\textbf{Esempio di famiglia pairwise independent}: \(p\) primo, \(\mathcal{K} = \mathbb{Z}_p \times \mathbb{Z}_p\), \(\mathcal{M} = \mathcal{T} = \mathbb{Z}_p\), \(h_{a,b}(m) = a \cdot m + b \bmod p\). Per \(m \ne m'\) e una qualunque coppia \((\tau, \tau')\), il sistema \(a m + b = \tau\), \(a m' + b = \tau'\) ha una sola soluzione \((a, b)\): sottraendo si ottiene \(a(m - m') = \tau - \tau'\), e \(m - m'\) è invertibile modulo \(p\). Quindi ogni coppia esce con probabilità \(1/p^2\).

\end{approfondimento}

\section{Riepilogo}

\begin{itemize}
\tightlist
\item
  \textbf{Shannon}: ogni schema perfettamente segreto ha \(|\mathcal{K}| \ge |\mathcal{M}|\), quindi chiavi lunghe almeno quanto i messaggi
\item
  Obiettivo d'ora in poi: una chiave di lunghezza indipendente dal messaggio
\item
  La segretezza perfetta si può definire con un gioco: nessun avversario distingue la cifratura di \(m_0\) da quella di \(m_1\)
\item
  \textbf{MAC}: Alice manda \((m, \tau)\) in chiaro con \(\tau = \text{Tag}(k, m)\), Bob ricalcola e confronta
\item
  \textbf{ε-sicurezza statistica one-time}: visto un tag valido, falsificarne un altro riesce al più con probabilità \(\varepsilon\); \(\varepsilon = 0\) è impossibile
\item
  L'OTP come MAC è completamente insicuro
\item
  Hash \textbf{pairwise independent} → MAC con \(\text{Tag}(k, m) = h_k(m)\)
\end{itemize}

% Generato da tools/notion2latex.py a partire dalla pagina Notion "lezione-3".

\chapter{Lezione 3 -- MAC da hash ed estrazione di casualità}

\section{Esempio di famiglia pairwise independent}

\[
h_{a,b}(m) = a \cdot m + b \bmod p
\]

\begin{itemize}
\tightlist
\item
  \(p\) è un \textbf{primo grande} (di \(\lambda\) bit)
\item
  Chiave: \((a, b) \in \mathbb{Z}_p^2 = \mathcal{K}\)
\item
  Messaggi e tag: \(m \in \mathbb{Z}_p = \mathcal{M} = \mathcal{T}\)
\end{itemize}

\subsection{Lemma}

\begin{inbreve}

Dimostriamo che \(a \cdot m + b \bmod p\) è pairwise independent, cioè che i tag di due messaggi diversi si comportano come due numeri casuali indipendenti.

\end{inbreve}

La famiglia \(\mathcal{H} = \{h_{a,b}\}\) qui sopra è \textbf{pairwise independent}.

\subsection{Dimostrazione}

\begin{inbreve}

Si contano le chiavi \((a, b)\) che danno una certa coppia di tag. È un sistema di 2 equazioni in 2 incognite, quindi c'è una sola chiave: ogni coppia di tag ha probabilità \(1/p^2\).

\end{inbreve}

Fissiamo \(m, m' \in \mathbb{Z}_p\) con \(m \ne m'\) e \(\tau, \tau' \in \mathbb{Z}_p\). Vogliamo calcolare

\[
\Pr_{a,b}[h_{a,b}(m) = \tau \wedge h_{a,b}(m') = \tau']
\]

Le due condizioni \(a m + b = \tau\) e \(a m' + b = \tau'\) formano un \textbf{sistema lineare} nelle incognite \(a, b\). In forma matriciale:

\[
\Pr_{a,b}\left[ \begin{pmatrix} m & 1 \\ m' & 1 \end{pmatrix} \begin{pmatrix} a \\ b \end{pmatrix} = \begin{pmatrix} \tau \\ \tau' \end{pmatrix} \right] = \Pr_{a,b}\left[ \begin{pmatrix} a \\ b \end{pmatrix} = \begin{pmatrix} m & 1 \\ m' & 1 \end{pmatrix}^{-1} \begin{pmatrix} \tau \\ \tau' \end{pmatrix} \right] = \frac{1}{p^2}
\]

\begin{itemize}
\tightlist
\item
  \textbf{L'inversa esiste} perché \(m \ne m'\): il determinante è \(m - m' \ne 0 \bmod p\), e modulo un primo ogni numero non nullo è invertibile.
\item
  Quindi c'è \textbf{una sola} coppia \((a, b)\) che soddisfa il sistema. Siccome \((a, b)\) è uniforme su \(p^2\) coppie, la probabilità di beccarla è \(1/p^2\). ∎
\end{itemize}

\textbf{Intuizione}: ogni coppia di tag \((\tau, \tau')\) corrisponde a esattamente una chiave, quindi la coppia di tag è uniforme su \(\mathbb{Z}_p^2\).

\section{MAC da hash pairwise independent}

\subsection{Teorema}

\begin{inbreve}

Se come tag usi una hash pairwise independent, Eve, anche dopo aver visto un messaggio con il suo tag, ha probabilità al massimo \(1/|\mathcal{T}|\) di produrre il tag giusto per un altro messaggio.
Eve conosce l'\textbf{insieme dei tag possibili} \(\mathcal{T}\) (fa parte dello schema, che è pubblico: es. ``i numeri da 0 a \(p-1\)''), ma non sa \textbf{quale} di questi è il tag giusto per il suo messaggio \(m'\), perché dipende dalla chiave. Può usare qualunque strategia (riciclare \(\tau\), fare calcoli, \ldots), ma nessuna funziona meglio che sceglierne uno a caso tra i \(|\mathcal{T}|\) possibili.

\end{inbreve}

Ogni famiglia pairwise independent \(\mathcal{H} = \{h_k : \mathcal{M} \to \mathcal{T}\}\) definisce un MAC \textbf{ε-statisticamente sicuro} con

\[
\varepsilon = \frac{1}{|\mathcal{T}|}
\]

\subsection{Dimostrazione}

Per ogni \(m \ne m'\) e ogni \(\tau, \tau'\), per definizione di probabilità condizionata:

\[
\Pr[\tau' = h_K(m') \mid \tau = h_K(m)] = \frac{\Pr[\tau' = h_K(m') \wedge \tau = h_K(m)]}{\Pr[h_K(m) = \tau]} = \frac{1/|\mathcal{T}|^2}{1/|\mathcal{T}|} = \frac{1}{|\mathcal{T}|}
\]

\begin{itemize}
\tightlist
\item
  \textbf{Numeratore}: per pairwise independence la coppia \((h_K(m), h_K(m'))\) è uniforme su \(\mathcal{T}^2\), quindi ogni coppia ha probabilità \(1/|\mathcal{T}|^2\).
\item
  \textbf{Denominatore}: sommando il numeratore su tutti i \(|\mathcal{T}|\) possibili \(\tau'\) si ottiene \(|\mathcal{T}| \cdot 1/|\mathcal{T}|^2 = 1/|\mathcal{T}|\). ∎
\end{itemize}

Per l'esercizio 1 della lezione 2, \(1/|\mathcal{T}|\) è il \textbf{miglior} \(\varepsilon\) possibile.

\begin{nota}

Con l'esempio \(h_{a,b}\): \(|\mathcal{T}| = p \approx 2^\lambda\), quindi \(\varepsilon \approx 2^{-\lambda}\), con una chiave di \(2\lambda\) bit.

\end{nota}

\subsection{Da dove vengono questi numeri}

\begin{inbreve}

Con una chiave di \(2\lambda\) bit, la costruzione \(h_{a,b}\) dà sicurezza \(2^{-\lambda}\). Sono quattro piccoli conti, a partire dal fatto che \textbf{λ è il numero di bit di} \(p\) (il prof sceglie un primo ``di λ bit'').

\end{inbreve}

\textbf{1. Perché} \(p \approx 2^\lambda\). Un numero di λ bit sta tra \(2^{\lambda-1}\) e \(2^\lambda\). Per esempio, i numeri da 3 bit vanno da \(100 = 4\) a \(111 = 7\), tutti vicini a \(2^3 = 8\). Quindi un primo di λ bit vale circa \(2^\lambda\).

\textbf{2. Perché} \(|\mathcal{T}| = p\). Il tag è \((a \cdot m + b) \bmod p\): il ``mod p'' si applica a \textbf{tutta} l'espressione (prima si calcola \(a \cdot m + b\), poi si prende il resto della divisione per \(p\)). Il resto di una divisione per \(p\) sta sempre tra \(0\) e \(p-1\), quindi i tag possibili sono \(0, 1, \dots, p-1\): sono \(p\) valori, e \(|\mathcal{T}| = p \approx 2^\lambda\).

\begin{esempio}

\textbf{Esempio} con \(p = 7\), \(a = 3\), \(m = 2\), \(b = 5\): \(a \cdot m + b = 11\), e \(11 \bmod 7 = 4\). Qualunque calcolo si faccia, il tag è sempre uno tra 0, 1, \ldots, 6.

\end{esempio}

\textbf{3. Perché} \(\varepsilon \approx 2^{-\lambda}\). Dal teorema, \(\varepsilon = 1/|\mathcal{T}| = 1/p \approx 1/2^\lambda = 2^{-\lambda}\).

\textbf{4. Perché la chiave è di} \(2\lambda\) \textbf{bit.} La chiave è la coppia \((a, b)\), e \(a\) e \(b\) sono numeri mod \(p\), cioè tra \(0\) e \(p-1\). Per scrivere in binario qualunque numero fino a \(p-1\) bastano λ bit, quindi \(a\) e \(b\) occupano λ bit ciascuno: in totale \(\lambda + \lambda = 2\lambda\) bit.

\begin{esempio}

\textbf{Esempio con λ = 3}:

\begin{itemize}
\tightlist
\item
  \(p = 7\) (in binario \(111\), cioè 3 bit), e \(7 \approx 2^3 = 8\)
\item
  tag possibili: \(0, \dots, 6\), cioè 7 valori
\item
  \(\varepsilon = 1/7 \approx 1/8 = 2^{-3}\)
\item
  \(a\) e \(b\) valgono da 0 a 6, cioè da \(000\) a \(110\) in binario: 3 bit ciascuno, chiave di \(3 + 3 = 6 = 2 \cdot 3\) bit
\end{itemize}

\end{esempio}

\begin{esempio}

\textbf{Esempio reale con λ = 128}: \(p\) è un primo di 128 bit, \(\varepsilon \approx 2^{-128}\) (Eve non falsifica praticamente mai), chiave di 256 bit.

\end{esempio}

\textbf{Il messaggio}: con una chiave di \(2\lambda\) bit si ottiene sicurezza \(2^{-\lambda}\). È proprio il minimo previsto dalla ``bad news'' qui sotto con \(t = 1\): \((1 + 1) \cdot \lambda = 2\lambda\).

\subsection{Bad news}

\begin{inbreve}

Con la sicurezza statistica la chiave del MAC deve crescere con il numero di messaggi da autenticare. Lo stesso problema dell'OTP.

\end{inbreve}

\textbf{Teorema}: ogni MAC \textbf{t-time} \(2^{-\lambda}\)-sicuro ha una chiave di dimensione (almeno) \((t+1) \cdot \lambda\).

\emph{t-time} = sicuro anche se Eve vede \(t\) coppie valide. Quindi, come per la segretezza perfetta, con la sicurezza statistica la chiave \textbf{cresce} con il numero di messaggi da autenticare. Con \(t = 1\) servono \(2\lambda\) bit, esattamente quelli di \(h_{a,b}\).

\section{Randomness extraction}

\textbf{Definizione}: la \emph{randomness extraction} (estrazione di casualità) è il processo che, partendo da una sorgente di casualità \textbf{imperfetta} (in cui alcuni valori escono più spesso di altri), produce bit \textbf{perfettamente casuali}, cioè uniformi. La funzione che lo fa si chiama \textbf{estrattore} (\(\text{Ext}\)).

\begin{esempio}

\textbf{Esempio}: hai solo una moneta truccata che fa testa il 70\% delle volte. Se la usi così com'è per creare una chiave, Eve sa che i bit valgono più spesso 1 che 0: la chiave non è uniforme e i teoremi visti (OTP, MAC) non valgono più. L'estrattore prende i lanci della moneta truccata e tira fuori bit come quelli di una moneta onesta.

Le sorgenti reali sono di questo tipo: un computer è deterministico, quindi la casualità la prende da misure fisiche (movimenti del mouse, tempi tra un tasto e l'altro, rumore elettrico), che sono imprevedibili solo in parte.

\end{esempio}

\subsection{Un estrattore in pratica (trucco di von Neumann)}

\emph{(Non negli appunti del prof: esempio per capire cosa fa un estrattore.)}

Hai una moneta piegata che fa testa il 60\% delle volte e ti servono bit perfettamente 50/50. Lanci la moneta \textbf{due volte} e guardi la coppia:

\begin{itemize}
\tightlist
\item
  \textbf{Testa poi Croce} → scrivi \textbf{0}
\item
  \textbf{Croce poi Testa} → scrivi \textbf{1}
\item
  \textbf{Testa Testa} oppure \textbf{Croce Croce} → butti via e rilanci
\end{itemize}

\textbf{Perché funziona}: \(\Pr[\text{TC}] = 0{,}6 \times 0{,}4 = 0{,}24\) e \(\Pr[\text{CT}] = 0{,}4 \times 0{,}6 = 0{,}24\). Le due probabilità sono identiche, quindi 0 e 1 escono esattamente al 50\%, anche se la moneta è truccata.

È questo che fa un estrattore: \textbf{entra roba sbilanciata, esce roba perfettamente bilanciata}. Il prezzo è che si butta via una parte dei lanci: si ottengono \textbf{meno bit} di quanti se ne sono raccolti.

\begin{nota}

Questo trucco funziona solo se i lanci sono indipendenti tra loro. L'estrattore del corso (hash + seme) funziona con fonti molto più generali, ma l'idea è la stessa.

\end{nota}

Come fanno Alice e Bob a ottenere \textbf{bit casuali}? È cruciale in crittografia:

\begin{itemize}
\tightlist
\item
  le \textbf{chiavi} devono essere \textbf{uniformi}
\item
  la casualità servirà anche più avanti per la sicurezza (es. sicurezza \textbf{CPA})
\end{itemize}

Si può pensare come un processo a più passi:

\begin{enumerate}
\def\labelenumi{\arabic{enumi}.}
\tightlist
\item
  \textbf{Randomness extraction}: partendo da una sorgente di casualità \textbf{imperfetta}, generare casualità ``pura'' (uniforme).
\item
  \textbf{Espansione}: espandere in modo \textbf{deterministico} quello che si è estratto, quanto serve (\textbf{pseudocasualità}).
\item
  In più, un modo per \textbf{rinfrescare} la casualità, da poter espandere di nuovo.
\end{enumerate}

Oggi ci occupiamo del passo 1.

\textbf{La domanda}: abbiamo una sorgente \(X\) e vogliamo progettare una funzione \textbf{deterministica} \(\text{Ext}\) tale che \(\text{Ext}(X)\) sia \textbf{uniforme}.

È \textbf{impossibile in generale}: se \(X\) è costante, anche \(\text{Ext}(X)\) lo è. Serve che \(X\) sia \textbf{imprevedibile} almeno in parte.

\subsection{Definizione: min-entropia}

\begin{inbreve}

La min-entropia misura quanto è imprevedibile una fonte, in bit. Min-entropia \(k\) vuol dire che la fonte è difficile da indovinare quanto \(k\) lanci di una moneta onesta, e quindi che se ne possono estrarre circa \(k\) bit casuali.

\end{inbreve}

La \textbf{min-entropia} di \(X\) è

\[
H_\infty(X) = -\log \max_x \Pr[X = x]
\]

(logaritmo in base 2)

\textbf{Intuizione}: \(\max_x \Pr[X = x]\) è la probabilità di indovinare \(X\) al primo colpo puntando sul valore più probabile. La min-entropia misura quanto è difficile questo ``miglior tentativo'': se \(H_\infty(X) = k\), il miglior tentativo riesce con probabilità \(2^{-k}\). Più min-entropia c'è, meglio è.

\begin{esempio}

\textbf{Esempi}:

\begin{itemize}
\item
  \(X = 0^n\) (costante): \(\max \Pr = 1\), quindi \(H_\infty(X) = -\log 1 = 0\)
\item
  \(X \equiv U_n\), cioè \(X \leftarrow \{0,1\}^n\) uniforme: \(\max \Pr = 2^{-n}\), quindi \(H_\infty(X) = n\)
\item
  \emph{(non negli appunti del prof)} \(X\) su 2 bit con \(\Pr[00] = 1/2\) e gli altri tre valori a \(1/6\): \(H_\infty(X) = -\log(1/2) = 1\). Due bit di lunghezza, ma solo 1 bit di min-entropia, perché la sorgente è sbilanciata verso 00.
\end{itemize}

\end{esempio}

\subsection{A cosa serve la min-entropia}

L'estrattore butta via una parte dei bit. La domanda naturale è: \textbf{quanti bit puliti si possono ottenere, al massimo?}

La min-entropia risponde proprio a questo: misura \textbf{quanti ``bit puliti'' ci sono nascosti dentro la fonte sporca}. Da una fonte con min-entropia \(k\) si possono estrarre al massimo circa \(k\) bit perfetti.

\begin{itemize}
\tightlist
\item
  \textbf{Moneta onesta}: ogni lancio vale 1 bit pulito. Min-entropia \textbf{1}.
\item
  \textbf{Moneta piegata (testa al 60\%)}: ogni lancio vale \textbf{meno} di un bit pulito, perché è un po' prevedibile. Min-entropia circa \textbf{0,74}.
\item
  \textbf{Moneta che fa sempre testa}: ogni lancio vale \textbf{0} bit puliti. Non c'è niente da estrarre.
\end{itemize}

Quindi la min-entropia è il \textbf{``contenuto di casualità vera''} della fonte, e dice quanta roba l'estrattore può tirarne fuori.

\subsection{Come si calcola}

La formula \(-\log_2(p)\) si può riscrivere in un modo più comodo:

\[
H_\infty(X) = \log_2\!\left(\frac{1}{p}\right) \qquad p = \text{probabilità del risultato più probabile}
\]

Cioè: si fa 1 diviso la probabilità, e poi ci si chiede ``2 elevato a quanto fa questo numero?'' (è proprio quello che fa \(\log_2\): ad esempio \(\log_2 8 = 3\) perché \(2^3 = 8\)).

\begin{esempio}

\textbf{Esempi}:

\begin{itemize}
\tightlist
\item
  \textbf{Moneta onesta}, \(p = 1/2\): \(1/p = 2\), e \(2^1 = 2\), quindi min-entropia \textbf{1}.
\item
  \textbf{Due monete oneste}, \(p = 1/4\): \(1/p = 4\), e \(2^2 = 4\), quindi min-entropia \textbf{2}.
\item
  \textbf{Moneta piegata}, \(p = 0{,}6\): \(1/p \approx 1{,}67\). Siccome \(2^0 = 1\) e \(2^1 = 2\), la risposta sta tra 0 e 1; il conto preciso dà \textbf{0,74} (infatti \(2^{0{,}74} \approx 1{,}67\)).
\end{itemize}

\end{esempio}

\textbf{Con la calcolatrice} (che di solito ha solo \(\ln\)): \(\log_2(x) = \ln(x) / \ln(2)\). Per la moneta piegata: \(\ln(1{,}67) / \ln(2) \approx 0{,}511 / 0{,}693 \approx 0{,}74\).

\textbf{Tipicamente} si sa solo che \(H_\infty(X) \ge k\) per qualche \(k < n\), e si vuole estrarre quanta più casualità possibile.

\subsection{Teorema: estrattori deterministici impossibili}

\begin{inbreve}

Un estrattore fissato in anticipo può sempre essere fregato da una fonte costruita apposta contro di lui, anche se quella fonte è quasi perfettamente casuale.

\end{inbreve}

\begin{nota}

\textbf{Notazione: cosa è cosa}

\begin{itemize}
\tightlist
\item
  \(\{0,1\}\) = l'insieme dei due bit possibili, 0 oppure 1: cioè \textbf{un singolo bit}.
\item
  \(\{0,1\}^n\) = le \textbf{stringhe di} \(n\) \textbf{bit}. Es. \(\{0,1\}^3\) contiene 000, 001, \ldots, 111 (8 stringhe). L'esponente dice quanti bit; senza esponente è uno solo.
\item
  \(\to\) = ``restituisce'': a sinistra cosa entra nella funzione, a destra cosa esce.
\item
  \(\times\) = ``e'': la funzione prende \textbf{due} input insieme.
\end{itemize}

Quindi \(\text{Ext}: \{0,1\}^n \to \{0,1\}\) si legge: \emph{``Ext prende una stringa di* \(n\) \emph{bit (quella della fonte) e restituisce un solo bit''}. Es. con \(n = 3\) e Ext =''prendi il primo bit'': \(\text{Ext}(101) = 1\).
Nella versione con il seme, \(\text{Ext}: \{0,1\}^d \times \{0,1\}^n \to \{0,1\}^\ell\) si legge: }''entra un seme di* \(d\) \emph{bit e una stringa della fonte di} \(n\) \emph{bit, escono} \(\ell\) \emph{bit puliti''}.

\end{nota}

\textbf{Perché un solo bit in output}: si sceglie l'obiettivo più facile possibile. Se non si riesce a estrarre nemmeno 1 bit pulito, a maggior ragione non si riesce con di più: così il risultato negativo è il più forte possibile.

Non esiste nessun \(\text{Ext}: \{0,1\}^n \to \{0,1\}\) che funzioni per tutte le sorgenti con \(H_\infty(X) \ge k\), \textbf{nemmeno} con \(k = n - 1\) (e nemmeno per estrarre un solo bit).

\subsection{Dimostrazione}

Prendiamo un qualunque \(\text{Ext}: \{0,1\}^n \to \{0,1\}\).

Sia \(b \in \{0,1\}\) il bit che \textbf{massimizza} \(|\text{Ext}^{-1}(b)|\), cioè il bit con più preimmagini.

\begin{verbatim}
 {0,1}^n
┌───────────────┐
│  x  x  x  x  │ ───→  0
│───────────────│
│  x  x  x     │ ───→  1
└───────────────┘
\end{verbatim}

Le due preimmagini \(\text{Ext}^{-1}(0)\) e \(\text{Ext}^{-1}(1)\) si spartiscono tutte le \(2^n\) stringhe, quindi la più grande ne contiene almeno metà:

\[
|\text{Ext}^{-1}(b)| \ge 2^{n-1}
\]

Sia \(X_{\text{BAD}}\) uniforme su \(\text{Ext}^{-1}(b)\).

\begin{itemize}
\tightlist
\item
  Ogni valore ha probabilità \(1/|\text{Ext}^{-1}(b)| \le 2^{-(n-1)}\), quindi \(H_\infty(X_{\text{BAD}}) \ge n - 1\)
\item
  \textbf{ma} \(\text{Ext}(X_{\text{BAD}}) = b\) è \textbf{costante}! ∎
\end{itemize}

\textbf{Intuizione}: qualunque estrattore fissato ha un ``punto debole''. Basta una sorgente cattiva che peschi solo dagli input che finiscono nello stesso output.

\subsection{Takeaway}

Nessun \(\text{Ext}\) deterministico può funzionare per \textbf{tutte} le sorgenti con min-entropia. Approcci possibili:

\begin{itemize}
\tightlist
\item
  \textbf{Più sorgenti indipendenti}: \emph{2-source extractors}.
\item
  \textbf{Strumenti crittografici} (es. funzioni hash): ma servono ipotesi forti su queste funzioni.
\item
  \textbf{Seeded extraction}: è quello che vediamo nel corso.
\end{itemize}

\section{Seeded extraction}

\begin{inbreve}

Invece di fissare l'estrattore, lo si sceglie a caso tra tantissimi usando un \textbf{seme} pubblico. Così nessuna fonte può essere costruita apposta contro l'estrattore che verrà usato.

\end{inbreve}

\textbf{Idea}: invece di un solo estrattore, una \textbf{famiglia} di estrattori, indicizzata da un \textbf{seed} (seme):

\[
\text{Ext}: \underbrace{\{0,1\}^d}_{\text{seed}} \times \underbrace{\{0,1\}^n}_{\text{sorgente}} \to \underbrace{\{0,1\}^\ell}_{\text{output}}
\]

Il seed sarà \textbf{casuale}. Perché non è barare (usare casualità per produrre casualità)?

\begin{itemize}
\tightlist
\item
  Il seed è un \textbf{catalizzatore} corto: l'output può essere più lungo.
\item
  Il seed è \textbf{pubblico} (l'avversario lo conosce!) e \textbf{riutilizzabile}.
\end{itemize}

\subsection{Definizione: estrattore con seed}

\begin{inbreve}

Un \((k, \varepsilon)\)-estrattore trasforma \textbf{ogni} fonte con min-entropia almeno \(k\) in bit che, presi insieme al seme, sono quasi indistinguibili da bit perfettamente casuali (a meno di un errore \(\varepsilon\)).

\end{inbreve}

\(\text{Ext}: \{0,1\}^d \times \{0,1\}^n \to \{0,1\}^\ell\) è un \((k, \varepsilon)\)\textbf{-seeded extractor} se per ogni \(X\) con \(H_\infty(X) \ge k\):

\[
(S, \text{Ext}(S, X)) \approx_\varepsilon (S, U_\ell) \ \equiv U_{d+\ell}
\]

dove \(S \leftarrow U_d\) (seed uniforme).

\textbf{Significato di} \(\approx_\varepsilon\): ogni avversario \(A\), anche con \textbf{potenza illimitata}, distingue \((S, \text{Ext}(S, X))\) da \(U_{d+\ell}\) con probabilità (vantaggio) al più \(\varepsilon\).

\textbf{Nota}: il seed \(S\) compare \textbf{insieme} all'output. È questo che formalizza ``il seed è pubblico'': l'output deve sembrare uniforme anche a chi conosce il seed.

\subsection{Distanza statistica}

\begin{inbreve}

La distanza statistica è un numero tra 0 e 1 che dice quanto due distribuzioni sono diverse. È il massimo vantaggio con cui chiunque, anche con potenza infinita, riesce a distinguerle.

\end{inbreve}

\[
X \approx_\varepsilon Y \iff \text{SD}(X; Y) \le \varepsilon
\]

\[
\text{SD}(X; Y) = \frac{1}{2} \sum_z \Big| \Pr[X = z] - \Pr[Y = z] \Big|
\]

\begin{nota}

\textbf{Cosa sono le variabili}

\begin{itemize}
\tightlist
\item
  \(X\), \(Y\) = due \textbf{distribuzioni}, cioè due modi di generare valori a caso (es. due monete).
\item
  \(z\) = \textbf{un valore possibile} (per una moneta: 1 o 0).
\item
  \(\Pr[X = z]\) = la probabilità che \(X\) dia il valore \(z\); idem \(\Pr[Y = z]\).
\item
  \(|\dots - \dots|\) = di quanto differiscono le due probabilità su quel valore (sempre positivo).
\item
  \(\sum_z\) = ``somma su tutti i valori \(z\) possibili'': si fa il conto per ogni valore e si sommano i risultati.
\item
  \(\frac{1}{2}\) = si divide per 2, perché quello che una distribuzione ha in più su un valore lo ha in meno su un altro (le probabilità sommano sempre a 1): senza il ½ la stessa differenza verrebbe contata due volte. Così la SD sta sempre tra 0 e 1.
\item
  \(\text{SD}(X; Y)\) = la \textbf{distanza statistica}, il numero che esce dalla formula.
\item
  \(\varepsilon\) = l'errore massimo che si accetta (un numero piccolo).
\item
  \(X \approx_\varepsilon Y\) = ``X e Y sono quasi uguali, con errore al massimo \(\varepsilon\)''.
\item
  \(\iff\) = ``se e solo se'', cioè ``è la stessa cosa di''. La prima formula è quindi solo la definizione del simbolo \(\approx_\varepsilon\).
\end{itemize}

\end{nota}

\textbf{Intuizione}: misura quanto due distribuzioni sono diverse. Vale 0 se sono identiche, 1 se non hanno valori in comune.

\begin{esempio}

\textbf{Esempio} \emph{(non negli appunti del prof)}: moneta truccata \(X\) con \(\Pr[1] = 0{,}6\) contro moneta onesta \(Y\):
\(\text{SD} = \frac{1}{2}(|0{,}6 - 0{,}5| + |0{,}4 - 0{,}5|) = 0{,}1\)

\end{esempio}

\section{Leftover Hash Lemma (LHL)}

\begin{inbreve}

Il \textbf{Leftover Hash Lemma} è il teorema che risponde alla domanda ``come si costruisce un estrattore con seme?''. La risposta: basta una famiglia di hash pairwise independent. Il seme sceglie la hash, la si applica alla fonte e si ottengono bit quasi perfettamente casuali, purché la fonte abbia un po' più min-entropia dei bit che si vogliono estrarre. Il nome viene da \emph{leftover} (avanzo): la casualità che ``avanza'' in una fonte imperfetta si trasforma con una hash in bit puliti.

\end{inbreve}

\subsection{Teorema}

Sia \(\mathcal{H} = \{h_s : \{0,1\}^n \to \{0,1\}^\ell\}_{s \in \{0,1\}^d}\) \textbf{pairwise independent}. Allora

\[
\text{Ext}(s, x) = h_s(x)
\]

è un \((k, \varepsilon)\)-seeded extractor per

\[
k \ge \ell + \underbrace{2\log(1/\varepsilon) - 2}_{\text{min-entropy gap}}
\]

\textbf{Cosa vuol dire} \(\text{Ext}(s, x) = h_s(x)\): il seme \(s\) sceglie quale funzione hash usare, e la si applica al valore \(x\) della fonte. L'output della hash sono i bit puliti.

\begin{esempio}

\textbf{Esempio} con \(h(x) = a \cdot x + b \bmod 7\): il seme è la coppia \(s = (a, b)\), per esempio \((3, 5)\). Se la fonte dà \(x = 2\), l'output è \(3 \cdot 2 + 5 = 11 \bmod 7 = 4\). Nei MAC la coppia \((a, b)\) era la \textbf{chiave segreta}, qui è il \textbf{seme pubblico}: stesse funzioni, ruolo diverso.

\end{esempio}

\textbf{Intuizione}: per estrarre \(\ell\) bit quasi uniformi servono un po' più di \(\ell\) bit di min-entropia. Il ``prezzo'' è il \textbf{min-entropy gap}, che cresce se vogliamo \(\varepsilon\) più piccolo.

\begin{esempio}

\textbf{Esempio} \emph{(non negli appunti del prof)}: per estrarre \(\ell = 128\) bit con \(\varepsilon = 2^{-64}\) serve \(k \ge 128 + 128 - 2 = 254\) bit di min-entropia.

\end{esempio}

\subsection{Collision probability e lemma}

\begin{inbreve}

La collision probability è la probabilità che due estrazioni indipendenti diano lo stesso valore. L'uniforme è la distribuzione che collide di meno, e il lemma dice che se una distribuzione collide quasi quanto l'uniforme, allora è quasi uniforme.

\end{inbreve}

La dimostrazione usa la \textbf{probabilità di collisione} di una variabile aleatoria \(Y\) su un insieme \(\mathcal{Y}\):

\[
\text{Col}(Y) = \sum_{y \in \mathcal{Y}} \Pr[Y = y]^2 = \Pr[Y = Y']
\]

dove \(Y'\) è una \textbf{copia indipendente} di \(Y\): è la probabilità che due estrazioni indipendenti diano lo stesso valore. La distribuzione uniforme ha la collision probability più bassa possibile, \(1/|\mathcal{Y}|\).

\textbf{Lemma}: se

\[
\text{Col}(Y) \le \frac{1}{|\mathcal{Y}|}(1 + 4\varepsilon^2)
\]

allora \(\text{SD}(Y; U) \le \varepsilon\), con \(U\) uniforme su \(\mathcal{Y}\).

\textbf{Intuizione}: se la collision probability è quasi quella dell'uniforme, allora \(Y\) è quasi uniforme.

\begin{approfondimento}{Dimostrazione del lemma (il prof la fa all'inizio della Lezione 4: lì
c'è la versione completa, passo per passo)}

Sia \(N = |\mathcal{Y}|\) e \(p_y = \Pr[Y = y]\). Per Cauchy-Schwarz (\(\sum |a_y| \le \sqrt{N \sum a_y^2}\)):

\(\text{SD}(Y; U) = \frac{1}{2}\sum_y |p_y - 1/N| \le \frac{1}{2}\sqrt{N \sum_y (p_y - 1/N)^2}\)

Sviluppando il quadrato: \(\sum_y (p_y - 1/N)^2 = \sum_y p_y^2 - \frac{2}{N}\sum_y p_y + \frac{1}{N} = \text{Col}(Y) - \frac{1}{N} \le \frac{4\varepsilon^2}{N}\)

Quindi \(\text{SD}(Y; U) \le \frac{1}{2}\sqrt{N \cdot 4\varepsilon^2 / N} = \varepsilon\). ∎

\end{approfondimento}

\subsection{Dimostrazione del LHL}

\begin{inbreve}

Si contano le collisioni dell'output. Ci sono quelle ``per caso'', che ci sono anche nell'uniforme, e quelle dovute alla fonte che ripete lo stesso valore, che sono poche grazie alla min-entropia. Quindi l'output collide quasi quanto l'uniforme e, per il lemma, è quasi uniforme.

\end{inbreve}

Poniamo \(Y = (S, h(S, X))\), e sia \(Y' = (S', h(S', X'))\) una copia indipendente. Calcoliamo la collision probability:

\[
\begin{aligned}
\text{Col}(Y) &= \Pr[Y = Y'] \\
&= \Pr[S = S' \wedge h(S, X) = h(S', X')] \\
&= \Pr[S = S' \wedge h(S, X) = h(S, X')] \\
&= \underbrace{\Pr[S = S']}_{\text{Col}(S) \,=\, 2^{-d}} \cdot \Pr[h(S, X) = h(S, X')]
\end{aligned}
\]

\begin{itemize}
\tightlist
\item
  \textbf{3° passaggio}: se \(S = S'\), possiamo sostituire \(S'\) con \(S\).
\item
  \textbf{4° passaggio}: i seed sono indipendenti dalle sorgenti, quindi le due probabilità si separano. \(S\) è uniforme su \(2^d\) valori, quindi \(\text{Col}(S) = 2^{-d}\).
\end{itemize}

Ora spezziamo in due casi, \(X = X'\) oppure \(X \ne X'\):

\[
\begin{aligned}
&= 2^{-d} \Big( \Pr[X = X'] \cdot 1 + \Pr[h(S, X) = h(S, X') \wedge X \ne X'] \Big) \\
&\le 2^{-d} \Big( 2^{-k} + \Pr[h(S, X) = h(S, X') \wedge X \ne X'] \Big) \\
&\le 2^{-d} \Big( 2^{-k} + 2^{-\ell} \Big)
\end{aligned}
\]

\begin{itemize}
\tightlist
\item
  \textbf{Primo termine}: se \(X = X'\) gli hash coincidono sicuramente (probabilità 1). E \(\Pr[X = X'] = \sum_x \Pr[X = x]^2 \le \max_x \Pr[X = x] \cdot \sum_x \Pr[X = x] \le 2^{-k}\), perché \(H_\infty(X) \ge k\).
\item
  \textbf{Secondo termine}: per input diversi \(x \ne x'\), per pairwise independence \(\Pr_S[h_S(x) = h_S(x')] = 2^\ell \cdot 2^{-2\ell} = 2^{-\ell}\) (le coppie \((t, t)\) con valori uguali sono \(2^\ell\), ognuna con probabilità \(2^{-2\ell}\)).
\end{itemize}

Raccogliendo \(2^{-\ell}\):

\[
\begin{aligned}
&= \frac{1}{2^{d+\ell}} \Big( 2^{\ell - k} + 1 \Big) \\
&\le \frac{1}{|\mathcal{Y}|} \Big( 2^{2 - 2\log(1/\varepsilon)} + 1 \Big) \\
&= \frac{1}{|\mathcal{Y}|} \Big( 1 + 4\varepsilon^2 \Big)
\end{aligned}
\]

\begin{itemize}
\tightlist
\item
  \(|\mathcal{Y}| = 2^{d+\ell}\) perché \(Y\) è una stringa di \(d + \ell\) bit.
\item
  Dall'ipotesi \(k \ge \ell + 2\log(1/\varepsilon) - 2\) segue \(\ell - k \le 2 - 2\log(1/\varepsilon)\).
\item
  \(2^{2 - 2\log(1/\varepsilon)} = 4 \cdot (2^{\log(1/\varepsilon)})^{-2} = 4\varepsilon^2\).
\end{itemize}

Per il lemma: \(\text{SD}(Y; U) \le \varepsilon\), cioè \((S, h(S, X)) \approx_\varepsilon U_{d+\ell}\). ∎

\section{Riepilogo}

\begin{itemize}
\tightlist
\item
  \(h_{a,b}(m) = am + b \bmod p\) è pairwise independent: il sistema lineare ha una sola soluzione perché \(m \ne m'\)
\item
  Ogni famiglia pairwise independent dà un MAC ε-sicuro con \(\varepsilon = 1/|\mathcal{T}|\) (ottimo)
\item
  \textbf{Bad news}: un MAC t-time \(2^{-\lambda}\)-sicuro richiede chiavi di \((t+1)\lambda\) bit
\item
  \textbf{Min-entropia} \(H_\infty(X) = -\log \max_x \Pr[X = x]\): quanto è difficile indovinare \(X\) al primo colpo
\item
  Nessun estrattore \textbf{deterministico} funziona per tutte le sorgenti, nemmeno con \(k = n - 1\) e un solo bit di output
\item
  \textbf{Seeded extractor}: con un seed casuale, corto e \textbf{pubblico}, \((S, \text{Ext}(S, X)) \approx_\varepsilon U_{d+\ell}\)
\item
  \textbf{Distanza statistica} \(\text{SD}\): misura quanto due distribuzioni differiscono
\item
  \textbf{Leftover Hash Lemma}: hash pairwise independent sono estrattori, se \(k \ge \ell + 2\log(1/\varepsilon) - 2\)
\end{itemize}

% Generato da tools/notion2latex.py a partire dalla pagina Notion "lezione-4".

\chapter{Lezione 4 -- Sicurezza computazionale e OWF}

\emph{Segue gli appunti del prof. Le parti aggiunte per capire meglio sono segnate ``non negli appunti del prof''. PRG e indistinguibilità computazionale sono nella} Lezione 5

\section{Fine lezione 3: dimostrazione del lemma sulla collision probability}

\begin{inbreve}

Nella lezione 3, per dimostrare il Leftover Hash Lemma, avevamo usato un lemma senza dimostrarlo: \textbf{se una distribuzione collide quasi quanto l'uniforme, allora è quasi uniforme}. Il prof lo dimostra qui, usando la disuguaglianza di Cauchy-Schwarz.

\end{inbreve}

\subsection{Promemoria: cosa dice il lemma}

Sia \(Y\) una variabile aleatoria che prende valori in un insieme \(\mathcal{Y}\), e sia \(U\) uniforme su \(\mathcal{Y}\). Se

\[
\text{Col}(Y) \le \frac{1}{|\mathcal{Y}|}\,(1 + 4\varepsilon^2)
\]

allora \(\text{SD}(Y; U) \le \varepsilon\).

\begin{nota}

\textbf{Cosa sono le cose}

\begin{itemize}
\tightlist
\item
  \(\mathcal{Y}\) = l'insieme dei valori possibili; \(|\mathcal{Y}|\) = quanti sono.
\item
  \(U\) = la distribuzione uniforme su \(\mathcal{Y}\): ogni valore ha probabilità \(1/|\mathcal{Y}|\).
\item
  \(\text{Col}(Y) = \sum_{y} \Pr[Y = y]^2\) = probabilità di collisione: la probabilità che due estrazioni indipendenti di \(Y\) diano lo stesso valore. L'uniforme ha la più bassa possibile, \(1/|\mathcal{Y}|\).
\item
  \(\text{SD}(Y; U)\) = distanza statistica tra \(Y\) e l'uniforme (lezione 3).
\end{itemize}

\end{nota}

\subsection{Strumento: prodotto scalare e Cauchy-Schwarz}

\emph{(Definizioni non negli appunti del prof, servono per seguire la dimostrazione.)}

\textbf{Prodotto scalare} di due vettori \(\vec{u} = (u_1, \dots, u_N)\) e \(\vec{v} = (v_1, \dots, v_N)\): si moltiplicano le componenti a coppie e si somma.

\[
\langle \vec{u}, \vec{v} \rangle = u_1 v_1 + u_2 v_2 + \dots + u_N v_N
\]

\begin{esempio}

\textbf{Esempio}: \(\langle (1, 2), (3, 4) \rangle = 1 \cdot 3 + 2 \cdot 4 = 11\). In particolare \(\langle \vec{u}, \vec{u} \rangle = u_1^2 + \dots + u_N^2\).

\end{esempio}

\textbf{Disuguaglianza di Cauchy-Schwarz}: il prodotto scalare di due vettori non supera mai la radice del prodotto dei loro ``quadrati'':

\[
\langle \vec{u}, \vec{v} \rangle \le \sqrt{\langle \vec{u}, \vec{u} \rangle \cdot \langle \vec{v}, \vec{v} \rangle}
\]

\begin{esempio}

\textbf{Esempio}: con \((1, 2)\) e \((3, 4)\): a sinistra 11; a destra \(\sqrt{5 \cdot 25} = \sqrt{125} \approx 11{,}18\). Infatti \(11 \le 11{,}18\).

\end{esempio}

\subsection{Dimostrazione (del prof)}

\textbf{Passo 1: riscrivere la SD come prodotto scalare.} Per definizione:

\[
\text{SD}(Y; U) = \frac{1}{2} \sum_{y \in \mathcal{Y}} \Big| \Pr[Y = y] - \frac{1}{|\mathcal{Y}|} \Big|
\]

Il prof introduce due abbreviazioni, per ogni valore \(y\):

\begin{itemize}
\tightlist
\item
  \(q_y = \Pr[Y = y] - \frac{1}{|\mathcal{Y}|}\): di quanto la probabilità di \(y\) si scosta da quella uniforme (può essere positivo o negativo)
\item
  \(s_y = +1\) se \(q_y \ge 0\), \(s_y = -1\) altrimenti: il \textbf{segno} di \(q_y\)
\end{itemize}

Siccome un numero moltiplicato per il suo segno dà il suo valore assoluto (\(q_y \cdot s_y = |q_y|\)):

\[
\text{SD}(Y; U) = \frac{1}{2} \sum_{y \in \mathcal{Y}} q_y\, s_y = \frac{1}{2} \langle \vec{q}, \vec{s} \rangle
\]

dove \(\vec{q} = (q_y)_{y \in \mathcal{Y}}\) e \(\vec{s} = (s_y)_{y \in \mathcal{Y}}\) sono i vettori con una componente per ogni valore \(y\).

\textbf{Passo 2: Cauchy-Schwarz.}

\[
\text{SD}(Y; U) = \frac{1}{2} \langle \vec{q}, \vec{s} \rangle \le \frac{1}{2} \sqrt{\langle \vec{q}, \vec{q} \rangle \cdot \langle \vec{s}, \vec{s} \rangle} = \frac{1}{2} \sqrt{\sum_{y \in \mathcal{Y}} q_y^2 \cdot |\mathcal{Y}|}
\]

Infatti \(\langle \vec{s}, \vec{s} \rangle = \sum_y s_y^2 = |\mathcal{Y}|\), perché ogni \(s_y\) vale \(\pm 1\), quindi \(s_y^2 = 1\), e ce n'è uno per ogni valore.

\textbf{Passo 3: calcolare} \(\sum_y q_y^2\). Si sviluppa il quadrato \((a - b)^2 = a^2 + b^2 - 2ab\):

\[
\sum_{y} q_y^2 = \sum_{y} \Big( \Pr[Y = y] - \frac{1}{|\mathcal{Y}|} \Big)^2 = \sum_{y} \Big( \Pr[Y = y]^2 + \frac{1}{|\mathcal{Y}|^2} - 2\,\frac{\Pr[Y = y]}{|\mathcal{Y}|} \Big)
\]

Si sommano i tre pezzi separatamente:

\begin{itemize}
\tightlist
\item
  \(\sum_y \Pr[Y = y]^2 = \text{Col}(Y)\), per definizione
\item
  \(\sum_y \frac{1}{|\mathcal{Y}|^2} = |\mathcal{Y}| \cdot \frac{1}{|\mathcal{Y}|^2} = \frac{1}{|\mathcal{Y}|}\), perché è lo stesso numero sommato \(|\mathcal{Y}|\) volte
\item
  \(\sum_y 2\,\frac{\Pr[Y = y]}{|\mathcal{Y}|} = \frac{2}{|\mathcal{Y}|} \sum_y \Pr[Y = y] = \frac{2}{|\mathcal{Y}|}\), perché le probabilità sommano a 1
\end{itemize}

Quindi, usando l'ipotesi su \(\text{Col}(Y)\):

\[
\sum_{y} q_y^2 = \text{Col}(Y) + \frac{1}{|\mathcal{Y}|} - \frac{2}{|\mathcal{Y}|} \le \frac{1 + 4\varepsilon^2}{|\mathcal{Y}|} - \frac{1}{|\mathcal{Y}|} = \frac{4\varepsilon^2}{|\mathcal{Y}|}
\]

\textbf{Passo 4: rimettere insieme.}

\[
\text{SD}(Y; U) \le \frac{1}{2} \sqrt{\frac{4\varepsilon^2}{|\mathcal{Y}|} \cdot |\mathcal{Y}|} = \frac{1}{2} \sqrt{4\varepsilon^2} = \frac{1}{2} \cdot 2\varepsilon = \varepsilon
\]

∎

\textbf{Intuizione}: \(\sum_y q_y^2\) misura quanto \(Y\) si scosta dall'uniforme ``al quadrato'', ed è esattamente l'\textbf{eccesso di collisioni} \(\text{Col}(Y) - 1/|\mathcal{Y}|\). Se l'eccesso è piccolo, Cauchy-Schwarz garantisce che anche la distanza statistica è piccola.

\section{Sicurezza computazionale}

\begin{inbreve}

Finora la sicurezza valeva contro avversari con potenza di calcolo \textbf{infinita}, e il prezzo erano limiti pesantissimi (chiavi lunghe quanto i messaggi, chiavi che crescono con l'uso). Per avere crittografia \textbf{dimostrabilmente sicura ed efficiente} bisogna cambiare qualcosa: si fanno \textbf{due ipotesi} in più.

\end{inbreve}

I limiti visti finora:

\begin{itemize}
\tightlist
\item
  \textbf{Cifratura (Shannon)}: la chiave deve essere lunga almeno quanto il messaggio e va usata una volta sola.
\item
  \textbf{MAC}: per autenticare \(t\) messaggi servono chiavi di \((t+1)\lambda\) bit.
\item
  \textbf{Casualità}: da una fonte con min-entropia \(k\) non si estraggono più di circa \(k\) bit puliti.
\end{itemize}

Le due ipotesi del prof:

\begin{enumerate}
\def\labelenumi{\arabic{enumi}.}
\tightlist
\item
  \textbf{L'avversario gira su una macchina ``efficiente''}: non ha potenza infinita.
\item
  \textbf{Esistono problemi che sembrano ``difficili''}, cioè che nessuna macchina efficiente sa risolvere.
\end{enumerate}

\section{Ipotesi (i): l'avversario è PPT}

\begin{inbreve}

L'avversario è un programma (una macchina di Turing) che lavora in \textbf{tempo polinomiale} e può usare \textbf{bit casuali}. Si chiama \textbf{PPT}: \emph{Probabilistic Polynomial Time}.

\end{inbreve}

\textbf{Macchina di Turing (TM)} \emph{(definizione non negli appunti del prof)}: il modello matematico di un computer. In pratica: \textbf{un algoritmo}. Dire ``l'avversario è una TM'' vuol dire ``l'avversario è un programma''.

\textbf{Tempo polinomiale}: si conta il numero di passi in funzione della lunghezza dell'input \(n\). Il programma è efficiente se fa \(O(n^c)\) passi per qualche costante \(c \in \mathbb{N}\).

\begin{nota}

\textbf{Cosa vuol dire} \(O(n^c)\): ``al massimo una costante per \(n^c\)''. Esempi: \(n^2\), \(5n^3 + n\), \(n^{100}\) sono polinomiali, cioè efficienti. \(2^n\) \textbf{non} lo è: con \(n = 128\) sono \(2^{128}\) passi, impossibili in pratica. Quindi l'avversario può fare tutto tranne i calcoli esponenziali, come provare tutte le chiavi possibili.

\end{nota}

\textbf{In più}: l'avversario può generare una quantità (polinomiale) di \textbf{bit casuali}, cioè può ``tirare monete'' durante il calcolo.

\[
\text{Adv} \equiv \text{PPT TM} \quad (\text{Probabilistic Polynomial Time Turing Machine})
\]

\section{Ipotesi (ii): esistono problemi difficili}

\begin{inbreve}

Idealmente vorremmo basare la sicurezza su \(P \neq NP\). Nessuno sa farlo; quello che si sa fare è basarla su problemi specifici che si credono difficili, come la \textbf{fattorizzazione}.

\end{inbreve}

\subsection{P e NP}

\begin{inbreve}

\textbf{P} = problemi che si sanno \textbf{risolvere} velocemente. \textbf{NP} = problemi in cui una soluzione, se qualcuno te la dà, si sa \textbf{controllare} velocemente. Ogni problema in P sta anche in NP.

\end{inbreve}

\begin{itemize}
\tightlist
\item
  \(P\) = tutti i problemi risolvibili da una TM in tempo polinomiale.
\item
  \(NP\) = tutti i problemi tali che una soluzione si può \textbf{verificare} da una TM in tempo polinomiale.
\end{itemize}

Esempi del prof di problemi in NP:

\begin{itemize}
\tightlist
\item
  \textbf{FACTORING}: dato \(n = p \cdot q\), trovare \(p\) e \(q\). Se qualcuno ti dà \(p\) e \(q\), controlli subito che \(p \cdot q = n\); trovarli da solo invece è difficile.
\item
  \textbf{3COL}: dato un grafo \(G\), trovare una \textbf{3-colorazione} \(\varphi\) di \(G\). \emph{(Definizione non negli appunti del prof: un grafo è un insieme di punti collegati da linee; una 3-colorazione assegna a ogni punto uno di 3 colori in modo che due punti collegati non abbiano mai lo stesso colore.)} Data una colorazione, controllare che sia giusta è facile (si guarda ogni linea); trovarla è difficile.
\end{itemize}

\begin{esempio}

\textbf{Esempio intuitivo} \emph{(non negli appunti del prof)}: il sudoku. Controllare un sudoku compilato è facile, risolverlo da zero no.

\end{esempio}

\textbf{La domanda P = NP?} È ``tutto ciò che si \textbf{controlla} facilmente si può anche \textbf{trovare} facilmente?''. Nessuno lo sa, ma quasi tutti credono di no (\(P \neq NP\)). P e NP sono \textbf{insiemi di problemi}, non problemi.

\begin{approfondimento}{Approfondimento (non negli appunti del prof): perché le domande sì/no}

Formalmente P e NP parlano di problemi \textbf{sì/no}: ``3233 ha un divisore minore di 60?''. Se la risposta è sì, una \textbf{prova} (es. ``53'') si controlla in fretta. NP = domande in cui la prova del sì si controlla velocemente; P = domande a cui si risponde velocemente da soli. \(P = NP\) vorrebbe dire che non serve mai qualcuno che ti suggerisca la prova: c'è sempre un metodo veloce per arrivarci da solo.

\end{approfondimento}

\subsection{Il teorema ideale}

\[
\textbf{Teorema (ideale)}: \text{ se } P \neq NP, \text{ allora } \Pi \text{ è sicuro.}
\]

\textbf{Perché sarebbe bellissimo}: restano solo due possibilità:

\begin{enumerate}
\def\labelenumi{\arabic{enumi}.}
\tightlist
\item
  \(\Pi\) è sicuro, oppure
\item
  \(P = NP\), che quasi nessuno crede
\end{enumerate}

Rompere \(\Pi\) vorrebbe dire risolvere il problema più famoso dell'informatica.

\textbf{Bad news}: \textbf{nessuno sa dimostrare teoremi di questo tipo.}

\subsection{Quello che si sa fare}

\[
\textbf{Teorema (esempio reale)}: \text{ se FACTORING è difficile, allora } \Pi \text{ è sicuro.}
\]

\textbf{È ancora utile?} Sì, leggendolo al contrario (\textbf{contronominale}): \emph{se} \(\Pi\) \emph{non fosse sicuro, allora FACTORING non sarebbe difficile, cioè esisterebbe una PPT TM che fattorizza.} E siccome da decenni nessuno ha trovato un algoritmo veloce per fattorizzare, è molto improbabile che qualcuno rompa \(\Pi\).

\begin{nota}

\textbf{Contronominale} \emph{(definizione non negli appunti del prof)}: ``se A allora B'' è equivalente a ``se non B allora non A''. Es. ``se piove, la strada è bagnata'' equivale a ``se la strada non è bagnata, non piove''.

\end{nota}

\section{One-way function (OWF)}

\begin{inbreve}

Una \textbf{funzione one-way} è facile da calcolare ma difficile da invertire: dato \(x\) si calcola \(f(x)\) in un attimo, ma dato \(f(x)\) nessun algoritmo efficiente riesce a trovare un input che dia quel risultato. Per il prof è \textbf{``la madre di tutte le ipotesi''}: è l'\textbf{ipotesi minima} per fare crittografia interessante.

\end{inbreve}

\subsection{Spiegazione semplice: la moltiplicazione}

\emph{(Non negli appunti del prof.)}

\textbf{Andata.} Ti do due numeri primi, 61 e 53: moltiplicarli è immediato, \(61 \times 53 = 3233\).

\textbf{Ritorno.} Ti do solo \textbf{3233} e ti chiedo quali due numeri ho moltiplicato. Non c'è un trucco: devi provare a dividere per 2, 3, 5, 7, 11, \ldots{} finché trovi un divisore.

\textbf{Con numeri enormi} (primi da 300 cifre) moltiplicare resta istantaneo, tornare indietro richiederebbe più tempo dell'età dell'universo.

\begin{itemize}
\tightlist
\item
  \(x\) = l'input (i due primi), \(f\) = la funzione (moltiplicarli), \(f(x)\) = il risultato (3233)
\item
  \textbf{invertire} = dato \(f(x)\), trovare \(x\)
\end{itemize}

\begin{nota}

\textbf{Funzione che NON è one-way}: \(f(x) = x + 5\). Da 15 si torna a 10 togliendo 5, istantaneo.

\end{nota}

\textbf{Perché serve}: serve qualcosa che \textbf{tu} sai fare e \textbf{Eve} no. Chi conosce l'input (es. la chiave) calcola tutto facilmente; Eve, che vede solo il risultato, non riesce a tornare indietro.

\subsection{Definizione (del prof)}

Una funzione \textbf{deterministica} \(f : \mathcal{X} \to \mathcal{Y}\) è \textbf{one-way} se per ogni avversario PPT \(A\):

\[
\Pr\big[\text{GAME}^{\text{owf}}_{A,f}(\lambda) = 1\big] \le \text{negl}(\lambda)
\]

cioè la probabilità che \(A\) \textbf{vinca} il gioco qui sotto è \textbf{trascurabile}.

\begin{nota}

\textbf{Cosa sono le cose}

\begin{itemize}
\tightlist
\item
  \textbf{deterministica} = stesso input, stesso output sempre (niente casualità dentro \(f\)).
\item
  \(\mathcal{X}\) = insieme degli input, \(\mathcal{Y}\) = insieme degli output. Tipicamente \(f : \{0,1\}^\lambda \to \{0,1\}^\lambda\): stringhe di \(\lambda\) bit in, stringhe di \(\lambda\) bit out.
\item
  \(\lambda\) = il \textbf{parametro di sicurezza}: la lunghezza dell'input in bit. Più è grande, più il problema è difficile.
\item
  \(\text{GAME}^{\text{owf}}_{A,f}(\lambda)\) = il gioco tra l'avversario \(A\) e il challenger, per la funzione \(f\), con parametro \(\lambda\). Vale 1 se \(A\) vince.
\item
  \(\text{negl}(\lambda)\) = una funzione \textbf{trascurabile} (vedi sotto).
\end{itemize}

\end{nota}

\subsection{Il gioco della one-way function}

\begin{verbatim}
Avversario A                          Challenger C
                                      x ← X      (x scelto a caso)
      ←──────────── y ─────────────   y = f(x)
      ──────────── x' ─────────────→
                                      Output 1 (WIN) se f(x') = y
\end{verbatim}

\begin{enumerate}
\def\labelenumi{\arabic{enumi}.}
\tightlist
\item
  Il challenger sceglie \(x\) \textbf{a caso} in \(\mathcal{X}\) (la freccia \(x \leftarrow \mathcal{X}\) vuol dire ``scelto uniformemente a caso'').
\item
  Calcola \(y = f(x)\) e lo dà ad \(A\).
\item
  \(A\) risponde con un \(x'\).
\item
  \(A\) \textbf{vince} se \(f(x') = y\).
\end{enumerate}

\textbf{Dettagli importanti} \emph{(non negli appunti del prof)}:

\begin{itemize}
\tightlist
\item
  \textbf{Basta una qualunque controimmagine}: \(A\) vince anche se \(x' \neq x\), purché \(f(x') = y\).
\item
  \textbf{L'input è casuale}: la difficoltà è richiesta \textbf{in media}, su input casuali, non solo per qualche input particolare.
\end{itemize}

\subsection{Esempi di OWF (candidate)}

Nessuno ha mai \textbf{dimostrato} che le OWF esistono: sono un'\textbf{ipotesi}. Questi sono i candidati del prof, cioè problemi che si credono difficili:

\begin{itemize}
\tightlist
\item
  \textbf{FACTORING}: \(f(p, q) = p \cdot q\); invertire = fattorizzare.
\item
  \textbf{DISCRETE LOG} (logaritmo discreto): \(f(x) = g^x \bmod p\); invertire = trovare l'esponente \(x\) dato \(g^x \bmod p\).
\item
  \textbf{LATTICES} (reticoli) e \textbf{DECODING RANDOM CODES} (decodificare codici casuali): problemi di algebra lineare ``rumorosa''. \emph{(Non negli appunti del prof: sono importanti perché si crede che resistano anche ai computer quantistici.)}
\end{itemize}

\subsection{OWF e P vs NP}

\emph{(Non negli appunti del prof: approfondimento per capire perché le OWF sono un'ipotesi più forte di P ≠ NP.)}

\begin{inbreve}

Se le OWF esistono, allora \(P \neq NP\). Il viceversa non è garantito: anche se \(P \neq NP\) le OWF potrebbero non esistere.

\end{inbreve}

\begin{approfondimento}{Dimostrazione: OWF ⇒ P ≠ NP}

Si dimostra al contrario: \textbf{se P = NP, allora nessuna funzione è one-way}.

P e NP parlano di domande sì/no, ma a noi serve \textbf{trovare} \(x\). Si fa come nel gioco delle 20 domande: si scopre \(x\) \textbf{un bit alla volta}, con domande del tipo ``esiste un \(x'\) che inizia con \(p\) e tale che \(f(x') = y\)?''.

Questa domanda sta in NP: se qualcuno propone \(x'\), si controlla velocemente (si guarda come inizia e si calcola \(f(x')\)). Se \(P = NP\), c'è un algoritmo veloce che risponde.

\textbf{L'invertitore}: parte dal prefisso vuoto; a ogni passo chiede ``esiste una controimmagine che continua con 0?''; se sì aggiunge 0, altrimenti aggiunge 1. Resta sempre vero che esiste una controimmagine con quel prefisso, quindi dopo \(\lambda\) passi ha una controimmagine completa.

\textbf{Esempio} con \(x = 101\): inizia con 0? No → 1. Inizia con 10? Sì → 10. Inizia con 100? No → 101 ✓.

Sono \(\lambda\) domande veloci, quindi tempo polinomiale, e successo con probabilità 1: \(f\) non è one-way. ∎

\end{approfondimento}

\textbf{Perché P ≠ NP non basta}: \(P \neq NP\) garantisce solo che esiste \textbf{almeno qualche} caso difficile (caso peggiore). In crittografia la chiave è scelta \textbf{a caso}, quindi serve che il problema sia difficile \textbf{quasi sempre} (caso medio). È come un produttore che garantisce che almeno uno su 1000 lucchetti è impossibile da scassinare: se ne prendi uno a caso, quasi sicuramente è facile.

\section{Funzioni trascurabili (negligible)}

\begin{inbreve}

``piccolo'' in crittografia si definisce in modo \textbf{asintotico}: una funzione è \textbf{trascurabile} se diventa più piccola dell'inverso di \textbf{qualunque} polinomio. È il modo formale di dire ``probabilità così piccola che si può ignorare''.

\end{inbreve}

\subsection{Definizione (del prof)}

Una funzione \(\varepsilon(\lambda)\) è \textbf{trascurabile} (\emph{negligible}) in \(\lambda\) se

\[
\varepsilon(\lambda) = O\Big(\frac{1}{p(\lambda)}\Big) \quad \text{per } \textbf{ogni} \text{ polinomio } p(\lambda) = \text{poly}(\lambda)
\]

Si scrive \(\varepsilon(\lambda) = \text{negl}(\lambda)\).

\textbf{Come leggerla} \emph{(non negli appunti del prof)}: per ogni polinomio, anche enorme come \(\lambda^{100}\), da un certo punto in poi \(\varepsilon(\lambda)\) sta sotto \(1/p(\lambda)\) e ci resta. Basta \textbf{un solo} polinomio per cui non succede e la funzione non è trascurabile.

\subsection{Esempi}

\textbf{Esempio del prof}: \(2^{-\lambda}\) è trascurabile. Confronto con \(1/\lambda^{10}\):

\begin{longtable}[]{@{}llll@{}}
\toprule\noalign{}
λ & 2\^{}(−λ) & 1/λ\^{}10 & Chi è più piccolo? \\
\midrule\noalign{}
\endhead
\bottomrule\noalign{}
\endlastfoot
10 & ≈ 0,001 & 0,0000000001 & 1/λ\^{}10 \\
50 & ≈ 9 · 10\^{}(−16) & ≈ 1 · 10\^{}(−17) & 1/λ\^{}10 \\
100 & ≈ 8 · 10\^{}(−31) & 1 · 10\^{}(−20) & \textbf{2\^{}(−λ)} \\
200 & ≈ 6 · 10\^{}(−61) & ≈ 1 · 10\^{}(−23) & \textbf{2\^{}(−λ)} \\
\end{longtable}

All'inizio \(2^{-\lambda}\) è più grande, ma da \(\lambda = 59\) in poi passa sotto e non risale più: si dimezza a ogni passo, mentre \(1/\lambda^{10}\) scende sempre più lentamente. Con qualunque altro polinomio succede lo stesso, solo più tardi.

\textbf{Controesempio} \emph{(non negli appunti del prof)}: \(1/\lambda^{10}\) \textbf{non} è trascurabile. Con il polinomio \(\lambda^{11}\) dovrebbe finire sotto \(1/\lambda^{11}\), ma \(1/\lambda^{10}\) è sempre più grande (es. con \(\lambda = 10\): \(10^{-10} > 10^{-11}\)).

\textbf{Perché conta}: un attacco che riesce con probabilità \(1/\lambda^{10}\) si può ripetere \(\lambda^{10}\) volte (tempo ancora polinomiale) e alla fine riesce. Con probabilità \(2^{-\lambda}\), anche ripetendolo \(\lambda^{10}\) volte si arriva al massimo a \(\lambda^{10} \cdot 2^{-\lambda}\), ancora minuscolo.

\subsection{Perché è comoda (proprietà del prof)}

\[
\text{negl}(\lambda) + \text{negl}(\lambda) = \text{negl}(\lambda) \qquad \text{poly}(\lambda) \cdot \text{negl}(\lambda) = \text{negl}(\lambda)
\]

\begin{approfondimento}{Dimostrazione delle due proprietà (non negli appunti del prof)}

\textbf{Idea}: per battere \(1/\lambda^c\) si usa la definizione con un esponente \textbf{più grande} di \(c\), così c'è margine per assorbire il fattore 2 o il polinomio.

\textbf{Somma.} Fissato \(c\), si applica la definizione con \(c+1\): da un certo punto \(\varepsilon_1, \varepsilon_2 < 1/\lambda^{c+1}\). Allora per \(\lambda \ge 2\): \(\varepsilon_1 + \varepsilon_2 < 2/\lambda^{c+1} = (2/\lambda) \cdot 1/\lambda^c \le 1/\lambda^c\). ✓

\textbf{Prodotto.} Un polinomio da un certo punto sta sotto \(\lambda^d\) per qualche \(d\). Fissato \(c\), si applica la definizione con \(c + d\): \(\varepsilon < 1/\lambda^{c+d}\), quindi \(p(\lambda) \cdot \varepsilon < \lambda^d / \lambda^{c+d} = 1/\lambda^c\). ∎

\end{approfondimento}

\section{Conclusione: i mondi computazionali}

\begin{inbreve}

A seconda delle ipotesi che si fanno, si ottengono ``mondi'' diversi, cioè diverse cose che la crittografia riesce a costruire. Con le sole OWF si costruisce tutta la crittografia \textbf{simmetrica}; per quella \textbf{a chiave pubblica} sembra servire qualcosa in più.

\end{inbreve}

\begin{longtable}[]{@{}
  >{\raggedright\arraybackslash}p{(\columnwidth - 4\tabcolsep) * \real{0.3333}}
  >{\raggedright\arraybackslash}p{(\columnwidth - 4\tabcolsep) * \real{0.3333}}
  >{\raggedright\arraybackslash}p{(\columnwidth - 4\tabcolsep) * \real{0.3333}}@{}}
\toprule\noalign{}
\begin{minipage}[b]{\linewidth}\raggedright
Mondo
\end{minipage} & \begin{minipage}[b]{\linewidth}\raggedright
Ipotesi
\end{minipage} & \begin{minipage}[b]{\linewidth}\raggedright
Cosa si può costruire
\end{minipage} \\
\midrule\noalign{}
\endhead
\bottomrule\noalign{}
\endlastfoot
\textbf{Minicrypt} & Esistono le OWF & SKE, MAC, PRNG \\
\textbf{Cryptomania} & Ipotesi più forti delle OWF & PKE, scambio di chiavi (TLS) \\
\end{longtable}

\begin{nota}

\textbf{Le sigle} \emph{(definizioni non negli appunti del prof)}

\begin{itemize}
\tightlist
\item
  \textbf{SKE} (\emph{Symmetric Key Encryption}): cifratura a chiave simmetrica, Alice e Bob hanno la stessa chiave (come l'OTP).
\item
  \textbf{MAC}: codici di autenticazione dei messaggi (lezioni 2-3).
\item
  \textbf{PRNG} (\emph{Pseudo-Random Number Generator}): generatori di bit pseudocasuali, cioè che \emph{sembrano} casuali (vedi Lezione 5).
\item
  \textbf{PKE} (\emph{Public Key Encryption}): cifratura a chiave pubblica, una chiave pubblica per cifrare e una privata per decifrare (es. RSA).
\item
  \textbf{Scambio di chiavi} (\emph{key exchange}): due persone che non si sono mai viste concordano una chiave segreta parlando su un canale pubblico. È quello che fa \textbf{TLS} ogni volta che apri un sito in HTTPS.
\end{itemize}

\end{nota}

\textbf{La freccia sbarrata} nel disegno del prof: dalle sole OWF \textbf{non} si sa ottenere la crittografia a chiave pubblica. Minicrypt non porta automaticamente a Cryptomania.

\begin{approfondimento}{I cinque mondi di Impagliazzo (dalle dispense, non negli appunti del
prof)}

Russell Impagliazzo (1995) ha descritto cinque mondi possibili. Non sappiamo in quale viviamo; si spera in uno degli ultimi due.

\begin{itemize}
\tightlist
\item
  \textbf{Algorithmica}: \(P = NP\), niente problemi difficili, niente crittografia computazionale.
\item
  \textbf{Heuristica}: \(P \neq NP\), ma i problemi difficili sono facili \textbf{in media}: inutile per le chiavi casuali.
\item
  \textbf{Pessiland}: esistono problemi difficili in media, ma \textbf{non} le OWF. Il mondo peggiore.
\item
  \textbf{Minicrypt}: esistono le OWF, ma non la crittografia a chiave pubblica.
\item
  \textbf{Cryptomania}: esistono OWF \textbf{e} crittografia a chiave pubblica.
\end{itemize}

\end{approfondimento}

\section{Riepilogo}

\begin{itemize}
\tightlist
\item
  \textbf{Lemma della collision probability}: si dimostra scrivendo la SD come prodotto scalare e usando \textbf{Cauchy-Schwarz}; l'eccesso di collisioni controlla la distanza dall'uniforme
\item
  \textbf{Sicurezza computazionale}: due ipotesi in più, avversario efficiente e problemi difficili
\item
  \textbf{Avversario = PPT TM}: tempo polinomiale \(O(n^c)\) e bit casuali
\item
  \textbf{P} = problemi che si risolvono velocemente; \textbf{NP} = problemi in cui una soluzione si controlla velocemente (es. FACTORING, 3COL)
\item
  Ideale: ``se P ≠ NP allora Π sicuro'', ma nessuno sa farlo. Reale: ``se FACTORING è difficile allora Π sicuro'', e al contrario: rompere Π ⇒ fattorizzare
\item
  \textbf{OWF}: facile da calcolare, difficile da invertire; nel gioco l'avversario vince con probabilità \(\le \text{negl}(\lambda)\); è l'ipotesi minima della crittografia
\item
  \textbf{Trascurabile}: \(O(1/p(\lambda))\) per ogni polinomio; es. \(2^{-\lambda}\); negl + negl = negl, poly · negl = negl
\item
  Candidati OWF: FACTORING, DISCRETE LOG, LATTICES, DECODING RANDOM CODES
\item
  \textbf{Minicrypt} (OWF): SKE, MAC, PRNG. \textbf{Cryptomania}: PKE, scambio di chiavi (TLS)
\end{itemize}

% Generato da tools/notion2latex.py a partire dalla pagina Notion "lezione-5".

\chapter{Lezione 5 -- Pseudocasualità, PRG e hard-core bit}

\emph{Segue gli appunti del prof. Le parti aggiunte per capire meglio sono segnate ``non negli appunti del prof''. Usa i concetti della Lezione 4: avversario PPT, funzioni trascurabili, one-way function.}

\section{Pseudocasualità}

\begin{inbreve}

Una funzione \textbf{pseudocasuale} è una funzione \textbf{deterministica} che prende una stringa uniforme (corta) e restituisce \textbf{tanti bit} che ``sembrano'' casuali. ``Sembrano'' ha un significato preciso: \textbf{nessun avversario PPT riesce a distinguerli da bit davvero casuali}, se non con probabilità trascurabile.

\end{inbreve}

\textbf{Il collegamento con la lezione 3} \emph{(non negli appunti del prof)}: l'estrattore tira fuori pochi bit puliti da una fonte sporca; la funzione pseudocasuale li usa come seme e ne genera quanti ne servono. È il passo 2 (``espansione'') visto nella lezione 3.

\begin{verbatim}
fonte sporca → [Estrattore] → seme (pochi bit casuali) → [PRG] → tanti bit pseudocasuali
\end{verbatim}

\section{Indistinguibilità computazionale}

\begin{inbreve}

Due distribuzioni sono \textbf{computazionalmente indistinguibili} (\(\approx_c\)) se nessun avversario PPT riesce a capire da quale delle due viene un campione, se non con vantaggio trascurabile. È la versione ``per avversari efficienti'' della distanza statistica della lezione 3.

\end{inbreve}

\subsection{Ensemble di distribuzioni}

Con la sicurezza computazionale tutto dipende dal parametro \(\lambda\) (più è grande, più l'avversario fatica). Per questo non si confrontano due distribuzioni singole, ma due \textbf{famiglie} di distribuzioni, una per ogni valore di \(\lambda\):

\[
X = \{X_\lambda\}_{\lambda \in \mathbb{N}} \qquad Y = \{Y_\lambda\}_{\lambda \in \mathbb{N}}
\]

\begin{nota}

\textbf{Cosa vuol dire} \emph{(non negli appunti del prof)}: \(X_\lambda\) è un modo di generare stringhe casuali che dipende da \(\lambda\). Una famiglia così si chiama \textbf{ensemble}. Esempio: \(\{U_\lambda\}\), dove \(U_\lambda\) è la distribuzione \textbf{uniforme} su stringhe di \(\lambda\) bit (\(U_1\) su 1 bit, \(U_2\) su 2 bit, \ldots).

\end{nota}

\subsection{Definizione (del prof)}

\(X \approx_c Y\) se per ogni avversario PPT \(A\):

\[
\Big| \Pr[A(X_\lambda) = 1] - \Pr[A(Y_\lambda) = 1] \Big| \le \text{negl}(\lambda)
\]

dove \(\Pr[A(X_\lambda) = 1]\) è un'abbreviazione per \(\Pr[A(z) = 1 : z \leftarrow X_\lambda]\).

\begin{nota}

\textbf{Cosa sono le cose}

\begin{itemize}
\tightlist
\item
  \(A\) = un programma efficiente (il \textbf{distinguisher}) che riceve una stringa \(z\) e risponde 1 (``secondo me viene da \(Y\)'') o 0 (``secondo me viene da \(X\)''). È l'assaggiatore alla cieca.
\item
  \(z \leftarrow X_\lambda\) = ``\(z\) è estratta dalla distribuzione \(X_\lambda\)''.
\item
  \(\Pr[A(z) = 1 : z \leftarrow X_\lambda]\) = quanto spesso \(A\) risponde 1 quando la stringa viene da \(X\). I due punti si leggono ``nella situazione in cui''.
\item
  La differenza in valore assoluto è il \textbf{vantaggio} di \(A\): quanto cambia il suo comportamento a seconda della provenienza della stringa.
\item
  \(\le \text{negl}(\lambda)\) = il vantaggio deve essere trascurabile, e questo per \textbf{ogni} \(A\) PPT.
\end{itemize}

\end{nota}

\begin{esempio}

\textbf{Esempi di vantaggio} \emph{(non negli appunti del prof)}: \(A\) dice 1 il 50\% delle volte con \(X\) e il 50\% con \(Y\): vantaggio 0, non distingue. \(A\) dice 1 il 10\% con \(X\) e il 90\% con \(Y\): vantaggio 0,8, distingue benissimo.

\end{esempio}

\subsection{La stessa cosa come gioco}

\(\text{GAME}(\lambda, b)\), con \(b \in \{0, 1\}\):

\begin{verbatim}
Avversario A                          Challenger C
                                      se b = 0: z ← X_λ
                                      se b = 1: z ← Y_λ
      ←──────────── z ─────────────
      ──────── b' ∈ {0,1} ─────────→
\end{verbatim}

Il challenger estrae \(z\) da \(X_\lambda\) (se \(b = 0\)) o da \(Y_\lambda\) (se \(b = 1\)) e lo dà ad \(A\), che risponde con \(b'\). La definizione diventa:

\[
\{\text{GAME}(\lambda, 0)\} \approx_c \{\text{GAME}(\lambda, 1)\}
\]

cioè la risposta di \(A\) ha (quasi) la stessa distribuzione nei due giochi. È lo stesso schema di GAME(0)/GAME(1) della lezione 2.

\subsection{Esempio concreto: due macchine}

\emph{(Non negli appunti del prof.)}

\textbf{Lo scenario.} Ci sono due ``macchine'' che producono stringhe di bit: la macchina \(X\) (es. l'uscita di un generatore pseudocasuale) e la macchina \(Y\) (es. bit davvero casuali). Ad \(A\) viene data \textbf{una} stringa e deve dire da quale macchina viene: risponde 1 (``viene da \(Y\)'') o 0 (``viene da \(X\)''). La domanda è: \textbf{esiste un} \(A\) \textbf{che si comporta in modo diverso a seconda della macchina da cui arriva la stringa?}

\textbf{Caso 1: un generatore scarso.}

\begin{itemize}
\tightlist
\item
  Macchina \(X\): produce 2 bit, ma sempre \textbf{uguali} tra loro: 00 oppure 11, metà e metà.
\item
  Macchina \(Y\): produce 2 bit davvero casuali: 00, 01, 10, 11, ognuno al 25\%.
\item
  Regola di \(A\): \textbf{``rispondi 1 se i due bit sono diversi''}.
\end{itemize}

Quanto spesso \(A\) risponde 1?

\begin{itemize}
\tightlist
\item
  Se la stringa viene da \(X\): i bit sono sempre uguali, quindi \textbf{mai}: \(\Pr[A(z) = 1 : z \leftarrow X] = 0\)
\item
  Se la stringa viene da \(Y\): i bit sono diversi in 01 e 10, metà dei casi: \(\Pr[A(z) = 1 : z \leftarrow Y] = 1/2\)
\end{itemize}

\textbf{Vantaggio} \(= |0 - 1/2| = 1/2\). \(A\) si comporta in modo molto diverso nei due casi, quindi \textbf{distingue} le due macchine: questo generatore \textbf{non} è indistinguibile da bit casuali (\(X \not\approx_c Y\)).

\textbf{Caso 2: un generatore buono.} Ora \(X\) è un buon PRG. Qualunque regola efficiente \(A\) provi, risponde 1 \textbf{quasi con la stessa frequenza} nei due casi, ad esempio:

\begin{itemize}
\tightlist
\item
  \(\Pr[A \text{ risponde } 1 \text{ con } X] = 0{,}5000000001\)
\item
  \(\Pr[A \text{ risponde } 1 \text{ con } Y] = 0{,}5\)
\end{itemize}

Il vantaggio è \(0{,}0000000001\), trascurabile. Se questo vale per \textbf{ogni} \(A\) efficiente, allora \(X \approx_c Y\).

\textbf{La formula letta in italiano}: \emph{``per ogni programma efficiente} \(A\)\emph{, la frequenza con cui} \(A\) \emph{dice 1 sulle stringhe di} \(X\) \emph{e quella con cui dice 1 sulle stringhe di} \(Y\) \emph{differiscono di una quantità trascurabile''}. \(\lambda\) è la lunghezza delle stringhe: il vantaggio deve diventare trascurabile al crescere di \(\lambda\).

\textbf{Il gioco è la stessa cosa}: il challenger tira una moneta \(b\); se \(b = 0\) prende la stringa dalla macchina \(X\), se \(b = 1\) dalla macchina \(Y\), e la dà ad \(A\), che risponde \(b'\) (il suo tentativo di indovinare \(b\)). ``GAME(λ, 0) ≈c GAME(λ, 1)'' vuol dire che \(A\) dice 1 con quasi la stessa frequenza sia quando la stringa viene da \(X\) sia quando viene da \(Y\): è esattamente la formula. Nella Lezione 2 le due macchine erano ``cifra \(m_0\)'' e ``cifra \(m_1\)''.

\subsection{Statistica vs computazionale}

\emph{(Non negli appunti del prof.)}

\begin{longtable}[]{@{}
  >{\raggedright\arraybackslash}p{(\columnwidth - 4\tabcolsep) * \real{0.3333}}
  >{\raggedright\arraybackslash}p{(\columnwidth - 4\tabcolsep) * \real{0.3333}}
  >{\raggedright\arraybackslash}p{(\columnwidth - 4\tabcolsep) * \real{0.3333}}@{}}
\toprule\noalign{}
\begin{minipage}[b]{\linewidth}\raggedright
\end{minipage} & \begin{minipage}[b]{\linewidth}\raggedright
Statistica (lezione 3)
\end{minipage} & \begin{minipage}[b]{\linewidth}\raggedright
Computazionale
\end{minipage} \\
\midrule\noalign{}
\endhead
\bottomrule\noalign{}
\endlastfoot
Simbolo & \(X \approx_\varepsilon Y\) & \(X \approx_c Y\) \\
Chi prova a distinguere & Chiunque, anche con potenza infinita & Solo avversari PPT \\
Vantaggio ammesso & Al più un numero fissato ε & Una funzione trascurabile di λ \\
\end{longtable}

L'indistinguibilità computazionale è \textbf{più debole}: due distribuzioni possono essere molto diverse statisticamente, ma nessun avversario efficiente riesce ad accorgersene. È proprio questo che rende possibili i PRG.

\section{Pseudorandom generator (PRG)}

\begin{inbreve}

Un \textbf{PRG} è una funzione deterministica che prende un seme casuale di \(\lambda\) bit e restituisce \(\lambda + \ell\) bit, cioè \textbf{più} bit di quanti ne ha ricevuti, in modo che l'uscita sia computazionalmente indistinguibile da bit davvero casuali.

\end{inbreve}

\subsection{Definizione (del prof)}

Una funzione \(G : \{0,1\}^\lambda \to \{0,1\}^{\lambda + \ell}\), con \(\ell > 0\), è un \textbf{PRG (sicuro)} se

\[
\{X_\lambda\} = \{G(U_\lambda)\}_\lambda \ \approx_c \ \{U_{\lambda + \ell}\}_\lambda = \{Y_\lambda\}
\]

\begin{nota}

\textbf{Cosa sono le cose}

\begin{itemize}
\tightlist
\item
  \(U_\lambda\) = un seme di \(\lambda\) bit scelto \textbf{uniformemente a caso}.
\item
  \(G(U_\lambda)\) = la distribuzione che si ottiene scegliendo un seme a caso e applicandogli \(G\): escono \(\lambda + \ell\) bit.
\item
  \(U_{\lambda + \ell}\) = \(\lambda + \ell\) bit davvero casuali.
\item
  \(\ell\) = lo \textbf{stretch} (``allungamento''): quanti bit in più escono rispetto al seme.
\item
  In parole: nessun avversario PPT distingue l'uscita di \(G\) da bit davvero casuali.
\item
  \emph{(Non negli appunti del prof)}: si intende anche che \(G\) sia calcolabile in tempo polinomiale, altrimenti non servirebbe a niente.
\end{itemize}

\end{nota}

\textbf{Notare} \emph{(non negli appunti del prof)}: \(G\) è \textbf{deterministica} (stesso seme, stessa uscita). Tutta la casualità viene dal seme; \(G\) la ``stira'' su più bit.

\subsection{Esercizio: nessun PRG è sicuro contro attaccanti illimitati}

\begin{inbreve}

Un PRG può essere sicuro \textbf{solo} contro avversari efficienti. Un avversario con potenza infinita lo rompe sempre, anche se il PRG allunga di un solo bit (\(G : \{0,1\}^\lambda \to \{0,1\}^{\lambda+1}\)).

\end{inbreve}

\textbf{Soluzione del prof.} L'avversario \(A\), dato \(z \in \{0,1\}^{\lambda+1}\), cerca un seme \(s \in \{0,1\}^\lambda\) tale che \(z = G(s)\). Se lo trova risponde 1, altrimenti 0.

\begin{itemize}
\tightlist
\item
  \textbf{Se} \(z\) \textbf{viene dal PRG} (\(z = G(s)\) con \(s \leftarrow U_\lambda\)): il seme esiste di sicuro, quindi \(\Pr[A(z) = 1] = 1\).
\item
  \textbf{Se} \(z\) \textbf{è uniforme} (\(z \leftarrow U_{\lambda+1}\)): \(\Pr[A(z) = 1] \le 1/2\).
\end{itemize}

Quindi il vantaggio è almeno \(1 - 1/2 = 1/2\), per niente trascurabile.

\textbf{Perché} \(\le 1/2\) \emph{(non negli appunti del prof)}: i semi sono \(2^\lambda\), quindi \(G\) produce al massimo \(2^\lambda\) stringhe diverse. Le stringhe di \(\lambda + 1\) bit sono \(2^{\lambda+1}\), il doppio. Quindi una stringa uniforme cade tra le uscite di \(G\) al massimo nella metà dei casi.

\textbf{Perché non funziona contro un avversario PPT} \emph{(non negli appunti del prof)}: per cercare il seme \(A\) deve provare tutti i \(2^\lambda\) semi, cioè fare un numero \textbf{esponenziale} di passi. Un avversario PPT non può farlo. Un PRG è sicuro proprio perché nessun avversario efficiente sa fare di meglio.

\section{Come si costruisce un PRG?}

\begin{inbreve}

Se esistono le one-way function, esistono PRG che allungano di quanto si vuole (un numero polinomiale di bit). La dimostrazione completa è difficile; il prof ne vede i due ingredienti principali.

\end{inbreve}

\textbf{Teorema}: se esistono le OWF, esistono PRG con stretch \(\ell(\lambda) = \text{poly}(\lambda)\).

I due ingredienti:

\begin{enumerate}
\def\labelenumi{\arabic{enumi}.}
\tightlist
\item
  \textbf{Stretch amplification}: passare da \(\ell(\lambda) = 1\) a \(\ell(\lambda) = \text{poly}(\lambda)\), cioè da un PRG che allunga di 1 bit a uno che allunga di quanto si vuole.
\item
  \textbf{Hard-core bit}: per costruire il PRG che allunga di 1 bit, partendo da una OWF.
\end{enumerate}

\section{Ingrediente 1: stretch amplification}

\begin{inbreve}

Basta un PRG che aggiunge \textbf{un solo bit} per ottenerne uno che ne aggiunge quanti ne vogliamo. L'idea è applicare \(G\) più volte di fila, tenendo da parte un bit a ogni passo.

\end{inbreve}

\subsection{Teorema}

Se \(G : \{0,1\}^\lambda \to \{0,1\}^{\lambda+1}\) è un PRG, allora per ogni \(\ell(\lambda) = \text{poly}(\lambda)\) esiste un PRG \(G^\ell : \{0,1\}^\lambda \to \{0,1\}^{\lambda + \ell}\).

\begin{nota}

\textbf{Notazione} \emph{(non negli appunti del prof)}: \(G^\ell\) \textbf{non} è una potenza: è solo il nome del nuovo PRG, quello che allunga di \(\ell\) bit.

\end{nota}

\subsection{La costruzione (del prof)}

\begin{verbatim}
s0 ──→ [G] ──s1──→ [G] ──s2──→ ... ──s(ℓ-1)──→ [G] ──→ sℓ
        │           │                           │
        ↓           ↓                           ↓
        b1          b2                          bℓ
\end{verbatim}

\begin{itemize}
\tightlist
\item
  \(s_0 \leftarrow U_\lambda\): il seme iniziale, casuale
\item
  per ogni \(i \in [\ell]\) (cioè \(i = 1, \dots, \ell\)): \((s_i, b_i) = G(s_{i-1})\), con \(s_i \in \{0,1\}^\lambda\) e \(b_i \in \{0,1\}\)
\item
  \textbf{output}: \((b_1, \dots, b_\ell, s_\ell)\)
\end{itemize}

In parole: \(G\) prende \(\lambda\) bit e ne restituisce \(\lambda + 1\). Si spezza l'uscita in due: \(\lambda\) bit diventano il \textbf{nuovo seme} \(s_i\), l'ultimo bit \(b_i\) si \textbf{mette da parte}. Si ripete \(\ell\) volte e alla fine si restituiscono tutti i bit messi da parte più l'ultimo seme: \(\ell + \lambda\) bit.

\subsection{Spiegazione semplice con un esempio}

\emph{(Non negli appunti del prof.)}

\textbf{L'idea: riusare l'uscita come nuovo seme.} \(G\) prende 3 bit e ne restituisce 4. Li dividiamo: \textbf{3 bit} li \textbf{rimettiamo dentro} \(G\) come nuovo seme (sembrano casuali, quindi vanno benissimo come seme), \textbf{1 bit} lo \textbf{teniamo da parte} come ``guadagno''. E ripetiamo: a ogni giro guadagniamo un bit.

\begin{nota}

\textbf{Analogia}: un distributore che, se metti una moneta, ti ridà \textbf{la moneta + una caramella}. Rimetti la moneta, ti ridà moneta + caramella. Dopo 3 giri hai \textbf{3 caramelle} e ancora la moneta in mano.

\end{nota}

\textbf{Passo per passo} (valori di \(G\) inventati, solo per vedere il meccanismo), seme iniziale \textbf{101}:

\begin{itemize}
\tightlist
\item
  \textbf{Giro 1}: \(G(101) = 0110\) → tieni da parte l'ultimo bit \textbf{0}, rimetti dentro \textbf{011}
\item
  \textbf{Giro 2}: \(G(011) = 1101\) → tieni da parte \textbf{1}, rimetti dentro \textbf{110}
\item
  \textbf{Giro 3}: \(G(110) = 0001\) → tieni da parte \textbf{1}, rimane \textbf{000}
\end{itemize}

\textbf{Alla fine} si mettono insieme i bit messi da parte (0, 1, 1) e l'ultimo seme rimasto (000): \textbf{011 000} = 6 bit. Partiti da \textbf{3} bit, ne abbiamo ottenuti \textbf{6}. Con \(\ell\) giri se ne ottengono \(3 + \ell\).

\textbf{I nomi nella formula del prof}:

\begin{itemize}
\tightlist
\item
  \(s_0\) = il seme iniziale (101)
\item
  \(s_i\) = il seme che si rimette dentro al giro \(i\) (011, 110, 000)
\item
  \(b_i\) = il bit messo da parte al giro \(i\) (0, 1, 1)
\item
  \((s_i, b_i) = G(s_{i-1})\) = ``al giro \(i\), \(G\) applicato al seme precedente dà il nuovo seme e un bit''
\item
  output \((b_1, \dots, b_\ell, s_\ell)\) = tutti i bit messi da parte, più l'ultimo seme
\item
  \(G^\ell\) = il nome di questo nuovo generatore, quello che fa \(\ell\) giri
\end{itemize}

\subsection{Come dimostrarlo: per riduzione}

\begin{inbreve}

Una \textbf{riduzione} dimostra ``se qualcuno sa rompere A, allora sa rompere anche B'', costruendo un attaccante per B che usa quello per A come sottoprogramma. Se B è sicuro per ipotesi, allora anche A lo è.

\end{inbreve}

Il piano del prof: \textbf{supponiamo che esista un avversario PPT} \(A\) \textbf{che rompe} \(G^\ell\)\textbf{; costruiamo un avversario PPT} \(B\) \textbf{che rompe} \(G\)\textbf{.} Siccome \(G\) è un PRG per ipotesi, \(B\) non può esistere, quindi nemmeno \(A\).

\begin{nota}

\textbf{``Rompere''} \emph{(non negli appunti del prof)} = distinguere l'uscita dall'uniforme con vantaggio \textbf{non trascurabile}, cioè almeno \(1/p(\lambda)\) per qualche polinomio \(p\).

\end{nota}

\textbf{Il problema} (i ``???'' del prof):

\begin{verbatim}
A ←── z' ∈ {0,1}^(λ+ℓ) ─── B ←────── z ─────── C      z ← G(U_λ)  oppure  z ← U_(λ+1)
          ???
\end{verbatim}

\(B\) riceve dal challenger una stringa \(z\) di \(\lambda + 1\) bit, ma \(A\) sa attaccare solo stringhe di \(\lambda + \ell\) bit. Come fa \(B\) a usare \(A\)? Per risolvere questo serve una tecnica nuova: l'argomento ibrido.

\subsection{Argomento ibrido (hybrid argument)}

\begin{inbreve}

Per dimostrare che due distribuzioni ``lontane'' \(X\) e \(Y\) sono indistinguibili, si costruisce una \textbf{catena di distribuzioni intermedie} (gli \textbf{ibridi}) in cui ogni anello differisce dal successivo per un dettaglio piccolissimo. Se nessun avversario nota i singoli passi, non nota nemmeno la differenza tra i due estremi.

\end{inbreve}

\textbf{Definizione (del prof)}: per mostrare \(X \approx_c Y\) si trovano ibridi \(H_0, H_1, \dots, H_\ell\) tali che:

\begin{enumerate}
\def\labelenumi{\arabic{enumi}.}
\tightlist
\item
  \(H_0 \equiv X\) e \(H_\ell \equiv Y\) (gli estremi sono le due distribuzioni da confrontare)
\item
  \(H_0 \approx_c H_1\), \(H_1 \approx_c H_2\), \ldots, \(H_{\ell-1} \approx_c H_\ell\) (ogni passo è indistinguibile)
\end{enumerate}

Allora, per la \textbf{disuguaglianza triangolare}, \(H_0 \approx_c H_\ell\), \textbf{purché} \(\ell(\lambda) = \text{poly}(\lambda)\).

\begin{nota}

\textbf{Analogia} \emph{(non negli appunti del prof)}: cambi una foto un pixel alla volta, finché diventa un'altra foto. Se a ogni passo nessuno nota la differenza con la versione precedente, e i passi sono ``pochi'' (polinomiali), nessuno distingue nemmeno la prima foto dall'ultima.

\end{nota}

\begin{approfondimento}{Perché funziona e perché serve ℓ polinomiale (non negli appunti del
prof)}

Chiamiamo \(p_i = \Pr[A(z) = 1 : z \leftarrow H_i]\). La \textbf{disuguaglianza triangolare} dice che la distanza diretta tra due numeri non supera la somma delle distanze passando per numeri intermedi:

\(|p_0 - p_\ell| \le |p_0 - p_1| + |p_1 - p_2| + \dots + |p_{\ell-1} - p_\ell|\)

A sinistra c'è il vantaggio di \(A\) sugli estremi; a destra la somma dei vantaggi sui singoli passi. Ogni passo ha vantaggio trascurabile, e i passi sono \(\ell\): il totale è al più \(\ell \cdot \text{negl}(\lambda)\). Se \(\ell\) è polinomiale, poly · negl = negl (Lezione 4), quindi il totale è ancora trascurabile.

Se invece i passi fossero esponenziali (es. \(2^\lambda\)), \(2^\lambda \cdot 2^{-\lambda} = 1\): non sarebbe più trascurabile e l'argomento crollerebbe.

\end{approfondimento}

\subsection{Gli ibridi nel nostro caso}

Gli estremi sono:

\begin{itemize}
\tightlist
\item
  \(H_0 \equiv G^\ell(U_\lambda)\): tutto pseudocasuale
\item
  \(H_\ell \equiv U_{\ell + \lambda}\): tutto davvero casuale
\end{itemize}

L'ibrido \(H_i\) (definizione del prof):

\begin{itemize}
\tightlist
\item
  \(b_1, \dots, b_i \leftarrow \{0,1\}\): i primi \(i\) bit sono \textbf{davvero casuali}
\item
  \(s_i \leftarrow U_\lambda\): il seme al passo \(i\) è \textbf{davvero casuale}
\item
  \((b_{i+1}, \dots, b_\ell, s_\ell) = G^{\ell - i}(s_i)\): il resto si calcola con la costruzione, facendo i rimanenti \(\ell - i\) passi
\end{itemize}

Cioè: \textbf{si sostituisce un bit pseudocasuale alla volta con un bit davvero casuale}, partendo da sinistra.

\begin{esempio}

\textbf{Esempio con} \(\ell = 3\) \emph{(non negli appunti del prof)}:

\begin{itemize}
\tightlist
\item
  \(H_0\): \(s_0\) casuale, poi \(b_1, b_2, b_3, s_3\) tutti calcolati con \(G\) = \(G^3(U_\lambda)\)
\item
  \(H_1\): \(b_1\) casuale, \(s_1\) casuale, poi \(b_2, b_3, s_3\) calcolati con \(G\)
\item
  \(H_2\): \(b_1, b_2\) casuali, \(s_2\) casuale, poi \(b_3, s_3\) calcolati con \(G\)
\item
  \(H_3\): \(b_1, b_2, b_3\) casuali e \(s_3\) casuale = \(U_{3+\lambda}\)
\end{itemize}

\end{esempio}

\textbf{Cosa cambia tra} \(H_0\) \textbf{e} \(H_1\) \emph{(non negli appunti del prof)}: solo il primo passo. In \(H_0\) la coppia \((s_1, b_1) = G(s_0)\) è \textbf{pseudocasuale}; in \(H_1\) è \textbf{davvero casuale}. Tutto il resto si calcola allo stesso modo a partire da \(s_1\).

\subsection{\texorpdfstring{La riduzione: \(H_0 \approx_c H_1\)}{La riduzione: H\_0 \textbackslash approx\_c H\_1}}

Ora la riduzione funziona. Il prof la fa concretamente per \(H_0 \approx_c H_1\).

\textbf{Supponiamo per assurdo} che esista un PPT \(A\) che distingue \(H_0\) da \(H_1\). Costruiamo \(B\) che rompe \(G\):

\begin{verbatim}
A ←── (b1, b2, ..., bℓ, sℓ) ─── B ←──── z ∈ {0,1}^(λ+1) ──── C     z ← G(U_λ)  oppure  z ← U_(λ+1)
  ──────────── b' ────────────→   ────────── b' ─────────→
\end{verbatim}

\(B\), ricevuto \(z \in \{0,1\}^{\lambda+1}\):

\begin{enumerate}
\def\labelenumi{\arabic{enumi}.}
\tightlist
\item
  spezza \(z = (b_1, s_1)\), con \(b_1 \in \{0,1\}\) e \(s_1 \in \{0,1\}^\lambda\)
\item
  calcola il resto con la costruzione: \((b_2, \dots, b_\ell, s_\ell) = G^{\ell-1}(s_1)\)
\item
  dà ad \(A\) la stringa \((b_1, b_2, \dots, b_\ell, s_\ell)\)
\item
  quando \(A\) risponde \(b'\), \(B\) risponde anche lui \(b'\) al challenger
\end{enumerate}

\textbf{Osservazione chiave (del prof)}:

\begin{itemize}
\tightlist
\item
  se \(z \leftarrow G(U_\lambda)\), la coppia \((b_1, s_1)\) è un'uscita di \(G\), quindi la stringa data ad \(A\) è distribuita come \(H_0\)
\item
  se \(z \leftarrow U_{\lambda+1}\), la coppia \((b_1, s_1)\) è davvero casuale, quindi la stringa è distribuita come \(H_1\)
\end{itemize}

\textbf{Per ipotesi}, \(A\) distingue \(H_0\) da \(H_1\):

\[
\Big| \Pr[A(b_1, \dots, b_\ell, s_\ell) = 1 : (b_1, \dots, s_\ell) \leftarrow \underbrace{G^\ell(U_\lambda)}_{= H_0}] - \Pr[A(b_1, \dots, b_\ell, s_\ell) = 1 : (b_1, \dots, s_\ell) \leftarrow H_1] \Big| \ge \frac{1}{p(\lambda)}
\]

per qualche polinomio \(p(\lambda)\). E per l'osservazione chiave:

\[
\Pr[B(z) = 1 : z \leftarrow G(U_\lambda)] = \Pr[A(b_1, \dots, b_\ell, s_\ell) = 1 : (b_1, \dots, s_\ell) \leftarrow H_0]
\]

\[
\Pr[B(z) = 1 : z \leftarrow U_{\lambda+1}] = \Pr[A(b_1, \dots, b_\ell, s_\ell) = 1 : (b_1, \dots, s_\ell) \leftarrow H_1]
\]

Quindi \(B\) distingue \(G(U_\lambda)\) da \(U_{\lambda+1}\) con vantaggio \(\ge 1/p(\lambda)\), cioè \textbf{rompe} \(G\). Ma \(G\) è un PRG: \textbf{contraddizione}. Inoltre \(B\) è PPT, perché chiama \(G\) al più \(\ell\) volte (polinomiale) e poi \(A\).

Con una dimostrazione molto simile, \textbf{per ogni} \(i\)\textbf{:} \(H_i \approx_c H_{i+1}\), e per l'argomento ibrido \(G^\ell(U_\lambda) = H_0 \approx_c H_\ell = U_{\lambda+\ell}\). ∎

\begin{approfondimento}{Il caso generale H\_i ≈c H\_(i+1) (non negli appunti del prof)}

È uguale, ma \(B\) inserisce \(z\) al passo \(i+1\) invece che al passo 1. Ricevuto \(z\):

\begin{enumerate}
\def\labelenumi{\arabic{enumi}.}
\tightlist
\item
  sceglie \(b_1, \dots, b_i\) \textbf{davvero casuali}
\item
  spezza \(z = (b_{i+1}, s_{i+1})\)
\item
  calcola \((b_{i+2}, \dots, b_\ell, s_\ell) = G^{\ell-i-1}(s_{i+1})\)
\item
  dà ad \(A\) la stringa \((b_1, \dots, b_\ell, s_\ell)\) e ne inoltra la risposta
\end{enumerate}

Se \(z \leftarrow G(U_\lambda)\), il passo \(i+1\) è pseudocasuale (come se \(s_i\) fosse uniforme e \((s_{i+1}, b_{i+1}) = G(s_i)\)): la stringa è distribuita come \(H_i\). Se \(z\) è uniforme, il passo \(i+1\) è davvero casuale: la stringa è distribuita come \(H_{i+1}\). Stesso ragionamento di prima.

\end{approfondimento}

\section{Ingrediente 2: hard-core bit}

\textbf{Da dove si prende il primo} \(G : \{0,1\}^\lambda \to \{0,1\}^{\lambda+1}\)? In teoria si può costruire partendo da \textbf{qualunque} OWF, ma non è molto efficiente. Più avanti si vedranno costruzioni efficienti basate su \textbf{teoria dei numeri}, \textbf{reticoli} (lattices), ecc.

L'ingrediente principale sono gli \textbf{hard-core bit}.

\subsection{Hard-core bit}

\begin{inbreve}

Una OWF nasconde il suo input \(x\) ``nel complesso'', ma potrebbe lasciar trapelare qualche pezzo di \(x\). Un \textbf{hard-core bit} è un bit \(h(x)\), calcolato da \(x\), che resta \textbf{impossibile da indovinare} anche conoscendo \(f(x)\): sembra un bit completamente casuale.

\end{inbreve}

Sia \(f : \{0,1\}^n \to \{0,1\}^n\) una OWF. La domanda del prof: \textbf{possiamo dire che un particolare bit di} \(x\) \textbf{è difficile da calcolare?}

\textbf{Intuitivamente}: \(h : \{0,1\}^n \to \{0,1\}\) è \textbf{hard-core} per \(f\) se \(h(x)\) è difficile da calcolare conoscendo \(f(x)\).

\textbf{Formalmente (in alternativa)}:

\[
(f(x), h(x)) \approx_c (f(x), u) \qquad u \leftarrow \{0,1\}
\]

\begin{nota}

\textbf{Cosa sono le cose}

\begin{itemize}
\tightlist
\item
  \(x\) = l'input segreto, scelto a caso.
\item
  \(h\) = una funzione che da \(x\) calcola \textbf{un solo bit} (es. ``il primo bit di \(x\)'', o lo XOR di alcuni bit di \(x\)).
\item
  \((f(x), h(x))\) = l'output della OWF \textbf{insieme} al bit \(h(x)\).
\item
  \((f(x), u)\) = l'output della OWF insieme a un bit \(u\) \textbf{completamente casuale}, che non c'entra niente con \(x\).
\item
  In parole: anche conoscendo \(f(x)\), il bit \(h(x)\) è indistinguibile da un bit tirato a caso.
\end{itemize}

\end{nota}

\textbf{Perché serve questa nozione} \emph{(non negli appunti del prof)}: essere one-way vuol dire solo che non si riesce a ricavare \textbf{tutto} \(x\). Non garantisce che non si ricavi \textbf{qualche pezzo} di \(x\). Ad esempio se \(f\) è one-way, anche \(f'(x) = f(x) \| x_1\) (che scrive in chiaro il primo bit di \(x\)) è one-way, ma il primo bit di \(x\) ovviamente non è nascosto. Un hard-core bit è un bit che invece è \textbf{garantito} nascosto.

\subsection{Perché è utile: un PRG da una one-way permutation}

Supponiamo che \(f : \{0,1\}^n \to \{0,1\}^n\) sia una \textbf{one-way permutation (OWP)}.

\begin{nota}

\textbf{Permutazione} \emph{(definizione non negli appunti del prof)}: una funzione da \(\{0,1\}^n\) a \(\{0,1\}^n\) in cui \textbf{ogni output viene da esattamente un input} (una biiezione). In pratica ``rimescola'' le stringhe senza perderne nessuna. Quindi: \textbf{se} \(s\) \textbf{è uniforme, anche} \(f(s)\) \textbf{è uniforme}. Una OWP è una OWF che è anche una permutazione.

\end{nota}

Allora il prof definisce:

\[
G(s) = f(s) \,\|\, h(s) \qquad s \leftarrow U_n
\]

(\(\|\) = \textbf{concatenazione}: si attaccano le due stringhe una dopo l'altra. Es. \(f(s) = 111\) e \(h(s) = 1\) danno \(f(s) \| h(s) = 1111\). Quindi gli \(n\) bit di \(f(s)\) più il bit \(h(s)\) fanno \(n+1\) bit: con \(n = 128\), escono 129 bit.)

\textbf{Perché è un PRG}, passo per passo:

\[
G(s) = f(s) \,\|\, h(s) \ \approx_c \ f(s) \,\|\, b \ \equiv \ U_{n+1} \qquad b \leftarrow \{0,1\}
\]

\begin{enumerate}
\def\labelenumi{\arabic{enumi}.}
\tightlist
\item
  \(f(s) \| h(s) \approx_c f(s) \| b\): è esattamente la definizione di hard-core bit: \(h(s)\) non si distingue da un bit casuale \(b\), anche conoscendo \(f(s)\).
\item
  \(f(s) \| b \equiv U_{n+1}\): \(f(s)\) è uniforme su \(n\) bit (perché \(f\) è una permutazione e \(s\) è uniforme), e \(b\) è un bit uniforme indipendente. Insieme fanno \(n + 1\) bit uniformi.
\end{enumerate}

Quindi \(G : \{0,1\}^n \to \{0,1\}^{n+1}\) allunga di 1 bit ed è un PRG. Con l'ingrediente 1 si allunga poi di quanto si vuole.

\subsection{In generale}

\begin{itemize}
\tightlist
\item
  Si riesce a far funzionare la stessa idea partendo da una qualunque \textbf{OWF} (non serve che sia una permutazione), ma è più complicato.
\item
  \textbf{Ogni} OWF \(f\) ammette un hard-core bit \(h\). \emph{(Non negli appunti del prof: è il teorema di Goldreich-Levin; l'hard-core bit è lo XOR di un sottoinsieme casuale dei bit di} \(x\)\emph{.)}
\end{itemize}

\subsection{\texorpdfstring{Domanda: esiste un unico \(h\) che funziona per ogni \(f\)?}{Domanda: esiste un unico h che funziona per ogni f?}}

\textbf{No.} L'idea del prof: preso un qualunque \(h\), si costruisce una funzione \(f_{\text{BAD}}\) tale che

\begin{itemize}
\tightlist
\item
  \(h\) \textbf{non} è hard-core per \(f_{\text{BAD}}\)
\item
  ma \(f_{\text{BAD}}\) è comunque una OWF
\end{itemize}

Va bene se \(f_{\text{BAD}}\) usa al suo interno un'altra OWF ausiliaria \(f : \{0,1\}^n \to \{0,1\}^n\).

\begin{approfondimento}{Come si costruisce f\_BAD (non negli appunti del prof)}

Si prende \(f_{\text{BAD}}(x) = f(x) \,\|\, h(x)\), cioè si scrive \(h(x)\) \textbf{in chiaro} nell'output.

\begin{itemize}
\tightlist
\item
  \(h\) \textbf{non è hard-core per} \(f_{\text{BAD}}\): chi vede \(f_{\text{BAD}}(x)\) legge direttamente \(h(x)\) nell'ultimo bit. Altro che difficile da calcolare.
\item
  \(f_{\text{BAD}}\) \textbf{è comunque one-way}: se qualcuno sapesse invertire \(f_{\text{BAD}}\), cioè trovare \(x'\) con \(f(x') \| h(x') = f(x) \| h(x)\), avrebbe in particolare \(f(x') = f(x)\), e quindi saprebbe invertire \(f\). Ma \(f\) è one-way. (È una riduzione!)
\end{itemize}

Quindi nessun \(h\) fissato funziona per tutte le OWF: l'hard-core bit va scelto in base alla funzione. È lo stesso schema del ``nessun estrattore fisso funziona per tutte le fonti'' della lezione 3.

\end{approfondimento}

\section{Riepilogo}

\begin{itemize}
\tightlist
\item
  \textbf{Pseudocasuale} = bit generati in modo deterministico da un seme corto, che nessun avversario PPT distingue da bit casuali (se non con probabilità trascurabile)
\item
  \(X \approx_c Y\): per ogni PPT \(A\), \(|\Pr[A(X_\lambda) = 1] - \Pr[A(Y_\lambda) = 1]| \le \text{negl}(\lambda)\); equivalente al gioco GAME(λ, 0) ≈c GAME(λ, 1)
\item
  \textbf{PRG}: \(G : \{0,1\}^\lambda \to \{0,1\}^{\lambda+\ell}\) con \(G(U_\lambda) \approx_c U_{\lambda+\ell}\)
\item
  \textbf{Nessun PRG è sicuro contro avversari illimitati}: basta cercare il seme (vantaggio ≥ 1/2), ma richiede tempo esponenziale
\item
  \textbf{OWF ⇒ PRG} con stretch polinomiale; due ingredienti:
\item
  \textbf{Stretch amplification}: da stretch 1 a stretch poly, applicando \(G\) in catena; si dimostra per \textbf{riduzione} usando l'\textbf{argomento ibrido} (ibridi \(H_i\): i primi \(i\) bit davvero casuali)
\item
  \textbf{Hard-core bit}: \((f(x), h(x)) \approx_c (f(x), u)\); con una OWP, \(G(s) = f(s) \| h(s)\) è un PRG con stretch 1
\item
  Ogni OWF ha un hard-core bit, ma nessun \(h\) funziona per tutte le OWF (\(f_{\text{BAD}}\))
\end{itemize}


\appendix
% Generato da tools/notion2latex.py a partire dalla pagina Notion "ripasso-probabilita".

\chapter{Ripasso di probabilità di base}

\section{1. Le basi}

\textbf{Spazio campionario} \(\Omega\): l'insieme di tutti i risultati possibili di un esperimento casuale.

\textbf{Evento} \(A\): un sottoinsieme di \(\Omega\), cioè un gruppo di risultati che ci interessa.

\begin{esempio}

Esempio: lanci un dado. \(\Omega = \{1,2,3,4,5,6\}\). L'evento ``esce pari'' è \(A = \{2,4,6\}\).

\end{esempio}

\subsection{Assiomi}

\begin{itemize}
\tightlist
\item
  \(0 \le \Pr[A] \le 1\): una probabilità sta sempre tra 0 (impossibile) e 1 (certo)
\item
  \(\Pr[\Omega] = 1\): qualcosa succede di sicuro
\item
  se \(A\) e \(B\) sono \textbf{disgiunti} (non possono succedere insieme): \(\Pr[A \cup B] = \Pr[A] + \Pr[B]\)
\end{itemize}

\begin{esempio}

Esempio: dado. ``Esce 1'' e ``esce 2'' sono disgiunti, quindi \(\Pr[\text{1 o 2}] = 1/6 + 1/6 = 1/3\).

\end{esempio}

\subsection{Distribuzione uniforme}

Tutti i risultati hanno la stessa probabilità. Allora:

\[
\Pr[A] = \frac{|A|}{|\Omega|} = \frac{\text{casi favorevoli}}{\text{casi possibili}}
\]

\begin{esempio}

Esempio: pesco una carta da un mazzo da 52. \(\Pr[\text{asso}] = 4/52 = 1/13\).

\end{esempio}

\subsection{Complementare}

\[
\Pr[\neg A] = 1 - \Pr[A]
\]

Utilissimo quando ``almeno una volta'' è difficile da contare ma ``mai'' è facile.

\begin{esempio}

Esempio: lanci un dado 4 volte. Probabilità di almeno un 6?

Contare direttamente è scomodo (un 6, due 6, \ldots). Il complementare è ``nessun 6'': ogni lancio non dà 6 con prob. \(5/6\), quindi

\(\Pr[\text{almeno un 6}] = 1 - (5/6)^4 = 1 - 625/1296 \approx 0{,}518\)

\end{esempio}

\subsection{Unione di eventi non disgiunti}

\[
\Pr[A \cup B] = \Pr[A] + \Pr[B] - \Pr[A \cap B]
\]

Si sottrae l'intersezione perché altrimenti la conti due volte.

\begin{esempio}

Esempio: pesco una carta. Probabilità che sia di cuori \textbf{o} una figura (J, Q, K)?

\(\Pr[\text{cuori}] = 13/52\), \(\Pr[\text{figura}] = 12/52\), \(\Pr[\text{figura di cuori}] = 3/52\)

\(\Pr = 13/52 + 12/52 - 3/52 = 22/52 \approx 0{,}42\)

\end{esempio}

\section{2. Union bound}

\[
\Pr[A_1 \cup \dots \cup A_m] \le \sum_{i=1}^{m} \Pr[A_i]
\]

Vale \textbf{sempre}, anche se gli eventi non sono disgiunti: sommare le probabilità al massimo sovrastima (perché conti più volte le intersezioni). Si usa quando il valore esatto è difficile ma basta un limite superiore.

\begin{esempio}

Esempio: hai 3 lampadine, ognuna si brucia entro l'anno con probabilità 1\%. La probabilità che \textbf{almeno una} si bruci è al più \(3 \times 1\% = 3\%\).

Il valore esatto (se sono indipendenti) è \(1 - 0{,}99^3 \approx 2{,}97\%\): il bound è molto vicino.

\end{esempio}

\section{3. Probabilità condizionata e indipendenza}

\subsection{Probabilità condizionata}

Con \(\Pr[B] > 0\):

\[
\Pr[A \mid B] = \frac{\Pr[A \cap B]}{\Pr[B]}
\]

Si legge ``probabilità di \(A\) sapendo che è successo \(B\)''. In pratica: restringi il mondo ai soli casi in cui vale \(B\), e guardi quanti di quelli soddisfano anche \(A\).

\begin{esempio}

Esempio: lanci un dado e un amico ti dice ``è uscito un numero pari''. Probabilità che sia 6?

\(\Pr[6 \mid \text{pari}] = \dfrac{\Pr[6]}{\Pr[\text{pari}]} = \dfrac{1/6}{1/2} = \dfrac{1}{3}\)

Ha senso: i casi possibili sono ora solo \(\{2,4,6\}\), e uno su tre è il 6.

\end{esempio}

\subsection{Indipendenza}

\(A\) e \(B\) sono \textbf{indipendenti} se

\[
\Pr[A \cap B] = \Pr[A] \cdot \Pr[B]
\]

che equivale a \(\Pr[A \mid B] = \Pr[A]\): sapere che è successo \(B\) non cambia nulla su \(A\).

\begin{esempio}

Esempio sì: lanci due monete. ``Prima testa'' e ``seconda testa'' sono indipendenti: \(\Pr[\text{TT}] = 1/2 \cdot 1/2 = 1/4\).

Esempio no: peschi due carte \textbf{senza rimetterle}. \(\Pr[\text{1ª asso}] = 4/52\), ma \(\Pr[\text{2ª asso} \mid \text{1ª asso}] = 3/51\): la prima pescata cambia la seconda.

Esempio sorprendente: dado, \(A\) = ``pari'' \(= \{2,4,6\}\), \(B\) = ``al massimo 4'' \(= \{1,2,3,4\}\).

\(\Pr[A \cap B] = \Pr[\{2,4\}] = 2/6 = 1/3\) e \(\Pr[A]\Pr[B] = 1/2 \cdot 2/3 = 1/3\). Sono indipendenti anche se riguardano lo stesso dado!

\end{esempio}

\textbf{Attenzione}: disgiunti ≠ indipendenti. Due eventi disgiunti con probabilità positiva sono anzi \textbf{dipendenti}: se succede uno, l'altro è impossibile.

\subsection{Probabilità totale}

Se \(B_1, \dots, B_m\) partizionano \(\Omega\) (coprono tutti i casi e non si sovrappongono):

\[
\Pr[A] = \sum_{i} \Pr[A \mid B_i]\,\Pr[B_i]
\]

Cioè: spezzi il problema in casi, calcoli ogni caso e fai la media pesata.

\begin{esempio}

Esempio: hai due urne. Urna 1: 3 palline rosse e 2 blu. Urna 2: 1 rossa e 4 blu. Lanci una moneta per scegliere l'urna, poi peschi una pallina.

\(\Pr[\text{rossa}] = \Pr[\text{rossa} \mid U_1]\Pr[U_1] + \Pr[\text{rossa} \mid U_2]\Pr[U_2] = \tfrac{3}{5}\cdot\tfrac{1}{2} + \tfrac{1}{5}\cdot\tfrac{1}{2} = 0{,}4\)

\end{esempio}

\subsection{Teorema di Bayes}

\[
\Pr[A \mid B] = \frac{\Pr[B \mid A]\,\Pr[A]}{\Pr[B]}
\]

Serve a ``invertire'' la condizionata: conosci \(\Pr[B \mid A]\) ma ti interessa \(\Pr[A \mid B]\).

\begin{esempio}

Esempio (il classico test medico): una malattia colpisce l'1\% della popolazione. Il test è positivo nel 99\% dei malati, ma dà anche un 5\% di falsi positivi sui sani. Fai il test ed è positivo: che probabilità hai di essere malato?

Prima \(\Pr[\text{pos}]\) con la probabilità totale:

\(\Pr[\text{pos}] = 0{,}99 \cdot 0{,}01 + 0{,}05 \cdot 0{,}99 = 0{,}0099 + 0{,}0495 = 0{,}0594\)

Poi Bayes:

\(\Pr[\text{malato} \mid \text{pos}] = \dfrac{0{,}99 \cdot 0{,}01}{0{,}0594} \approx 0{,}167\)

Solo il 17\% circa! Il motivo: i sani sono tantissimi, quindi anche un 5\% di falsi positivi su di loro supera di molto i veri positivi.

\end{esempio}

\section{4. Variabili aleatorie e valore atteso}

\subsection{Variabile aleatoria}

Una \textbf{variabile aleatoria} \(X\) è una funzione \(\Omega \to \mathbb{R}\): associa un numero a ogni risultato dell'esperimento.

\begin{esempio}

Esempio: lanci due dadi, \(X\) = somma dei due. Se esce (3, 5), \(X = 8\). \(\Pr[X = 2] = 1/36\) (solo (1,1)), \(\Pr[X = 7] = 6/36\) (sei combinazioni).

\end{esempio}

\subsection{Valore atteso}

\[
\mathbb{E}[X] = \sum_x x \cdot \Pr[X = x]
\]

È la media dei valori pesata con le loro probabilità: il risultato ``medio'' se ripeti l'esperimento tantissime volte.

\begin{esempio}

Esempio: un dado. \(\mathbb{E}[X] = (1+2+3+4+5+6)/6 = 3{,}5\). Non esce mai 3,5, ma è la media a lungo andare.

Esempio (scommessa): paghi 1 € per giocare, se esce 6 vinci 5 €. Il guadagno netto \(X\) vale \(+4\) con prob. \(1/6\) e \(-1\) con prob. \(5/6\).

\(\mathbb{E}[X] = 4 \cdot \tfrac{1}{6} - 1 \cdot \tfrac{5}{6} = -\tfrac{1}{6} \approx -0{,}17\) €. In media perdi 17 centesimi a giocata.

\end{esempio}

\subsection{Linearità del valore atteso}

\[
\mathbb{E}[X + Y] = \mathbb{E}[X] + \mathbb{E}[Y] \qquad \mathbb{E}[aX] = a\,\mathbb{E}[X]
\]

Vale \textbf{sempre}, anche se \(X\) e \(Y\) non sono indipendenti.

\begin{esempio}

Esempio: somma di due dadi. Invece di fare la tabella 6×6: \(\mathbb{E}[X_1 + X_2] = 3{,}5 + 3{,}5 = 7\).

\end{esempio}

\subsection{Variabili indicatrici}

\(X_A = 1\) se l'evento \(A\) succede, \(0\) altrimenti. Allora \(\mathbb{E}[X_A] = \Pr[A]\).

Combinate con la linearità servono a \textbf{contare} quante volte succede qualcosa in media: scrivi il conteggio come somma di indicatrici.

\begin{esempio}

Esempio: lanci 10 monete. Numero medio di teste? \(X = X_1 + \dots + X_{10}\) dove \(X_i = 1\) se la moneta \(i\) dà testa. \(\mathbb{E}[X] = 10 \cdot 1/2 = 5\).

Esempio (guardaroba): \(n\) persone lasciano il cappello, che vengono restituiti a caso. Quante persone in media riavranno il proprio?

\(X_i = 1\) se la persona \(i\) riprende il suo: \(\Pr = 1/n\). Quindi \(\mathbb{E}[X] = n \cdot 1/n = 1\), \textbf{qualunque sia} \(n\). Le \(X_i\) non sono indipendenti, ma la linearità funziona lo stesso.

\end{esempio}

\subsection{Variabili indipendenti}

\(X\) e \(Y\) sono indipendenti se \(\Pr[X=x \wedge Y=y] = \Pr[X=x]\Pr[Y=y]\) per ogni \(x, y\). In tal caso vale (senza indipendenza in generale no):

\[
\mathbb{E}[XY] = \mathbb{E}[X]\,\mathbb{E}[Y]
\]

\begin{esempio}

Esempio: prodotto di due dadi lanciati separatamente: \(\mathbb{E}[X_1 X_2] = 3{,}5 \cdot 3{,}5 = 12{,}25\).

Controesempio: \(X\) = un dado, \(Y = X\) (lo stesso dado). \(\mathbb{E}[X \cdot X] = (1+4+9+16+25+36)/6 \approx 15{,}17 \ne 12{,}25\).

\end{esempio}


\end{document}
